%Version 3.1 December 2024 See section 11 of the User Manual for version history
%
%%%%%%%%%%%%%%%%%%%%%%%%%%%%%%%%%%%%%%%%%%%%%%%%%%%%%%%%%%%%%%%%%%%%%%
%%                                                                 %%
%% Please do not use \input{...} to include other tex files.       %% % Submit
%your LaTeX manuscript as one .tex document.              %%
%%                                                                 %%
%% All additional figures and files should be attached             %% %
%separately and not embedded in the \TeX\ document itself.       %%
%%                                                                 %%
%%%%%%%%%%%%%%%%%%%%%%%%%%%%%%%%%%%%%%%%%%%%%%%%%%%%%%%%%%%%%%%%%%%%%

%%\documentclass[referee,sn-basic]{sn-jnl}% referee option is meant for double
%line spacing

%%=======================================================%%
%% to print line numbers in the margin use lineno option %%
%%=======================================================%%

%%\documentclass[lineno,pdflatex,sn-basic]{sn-jnl}% Basic Springer Nature
%Reference Style/Chemistry Reference Style

%%=========================================================================================%%
%% the documentclass is set to pdflatex as default. You can delete it if not
%appropriate.  %%
%%=========================================================================================%%

%%\documentclass[sn-basic]{sn-jnl}% Basic Springer Nature Reference
%Style/Chemistry Reference Style

%%Note: the following reference styles support Namedate and Numbered
%referencing. By default the style follows the most common style. To switch
%between the options you can add or remove �Numbered� in the optional
%parenthesis. %The option is available for: sn-basic.bst, sn-chicago.bst%  
 
%%\documentclass[pdflatex,sn-nature]{sn-jnl}% Style for submissions to Nature
%Portfolio journals %\documentclass[pdflatex,sn-basic]{sn-jnl}% Basic Springer
%Nature Reference Style/Chemistry Reference Style
\documentclass[pdflatex,sn-mathphys-num]{sn-jnl}% Math and Physical Sciences
%% Numbered Reference Style %\documentclass[pdflatex,sn-mathphys-ay]{sn-jnl}%
%Math and Physical Sciences Author Year Reference Style
%%\documentclass[pdflatex,sn-aps]{sn-jnl}% American Physical Society (APS)
%Reference Style %\documentclass[pdflatex,sn-vancouver-num]{sn-jnl}% Vancouver
%Numbered Reference Style %\documentclass[pdflatex,sn-vancouver-ay]{sn-jnl}%
%Vancouver Author Year Reference Style %\documentclass[pdflatex,sn-apa]{sn-jnl}%
%APA Reference Style %\documentclass[pdflatex,sn-chicago]{sn-jnl}% Chicago-based
%Humanities Reference Style

%%%% Standard Packages %<additional latex packages if required can be included
%here>

\usepackage{graphicx}%
\usepackage{multirow}%
\usepackage{amsmath,amssymb,amsfonts}%
\usepackage{amsthm}%
\usepackage{mathrsfs}%
\usepackage[title]{appendix}%
\usepackage{xcolor}%
\usepackage{textcomp}%
\usepackage{manyfoot}%
\usepackage{booktabs}%
\usepackage{algorithm}%
\usepackage{algorithmicx}%
\usepackage{algpseudocode}%
\usepackage{listings}%
%%%%

%%%%%=============================================================================%%%%
%%%%  Remarks: This template is provided to aid authors with the preparation %%%
%of original research articles intended for submission to journals published %%%
%by Springer Nature. The guidance has been prepared in partnership with %%%
%production teams to conform to Springer Nature technical requirements. %%%
%Editorial and presentation requirements differ among journal portfolios and %%%
%research disciplines. You may find sections in this template are irrelevant %%%
%to your work and are empowered to omit any such section if allowed by the %%%
%journal you intend to submit to. The submission guidelines and policies %%%  of
%the journal take precedence. A detailed User Manual is available in the %%%
%template package for technical guidance.
%%%%%=============================================================================%%%%

%% as per the requirement new theorem styles can be included as shown below
\theoremstyle{thmstyleone}%
\newtheorem{theorem}{Theorem}%  meant for continuous numbers
%%\newtheorem{theorem}{Theorem}[section]% meant for sectionwise numbers %
%optional argument [theorem] produces theorem numbering sequence instead of
%independent numbers for Proposition
\newtheorem{proposition}[theorem]{Proposition}% 
%%\newtheorem{proposition}{Proposition}% to get separate numbers for theorem and
%proposition etc.

\theoremstyle{thmstyletwo}%
\newtheorem{example}{Example}%
\newtheorem{remark}{Remark}%

\theoremstyle{thmstylethree}%
\newtheorem{definition}{Definition}%

\raggedbottom
%%\unnumbered% uncomment this for unnumbered level heads

\begin{document}

\title[The Legend Challenge: Embedding Ethical Compliance into LLMs through Champion Narratives in the Gender Equality Domain]{The Legend Challenge: Embedding Ethical Compliance into LLMs through Champion Narratives in the Gender Equality Domain}

%%=============================================================%%
%% GivenName    -> \fnm{Joergen W.} % Particle -> \spfx{van der} -> surname
%prefix % FamilyName   -> \sur{Ploeg} % Suffix   -> \sfx{IV} %
%\author*[1,2]{\fnm{Joergen W.} \spfx{van der} \sur{Ploeg} %
%\sfx{IV}}\email{iauthor@gmail.com}
%%=============================================================%%

\author*[1]{\fnm{Fatmir} \sur{Bylyshi}}\email{fatmir.bylyshi@studenti.unimi.it}
\equalcont{These authors contributed equally to this work.}

\author*[1]{\fnm{Ouardi} \sur{Ilyass}}\email{ilyass.ouardi@studenti.unimi.it}
\equalcont{These authors contributed equally to this work.}

% \author[1,2]{\fnm{Third} \sur{Author}}\email{iiiauthor@gmail.com}
% \equalcont{These authors contributed equally to this work.}

\affil*[1]{\orgdiv{Department of Computer Science}, \orgname{University of
Milan}, \orgaddress{\street{Via Celoria, 20}, \city{Milan}, \postcode{20134},
\country{Italy}}}

%%==================================%%
%% Sample for unstructured abstract %%
%%==================================%%

% \abstract{Large companies operating under increasingly complex sustainability and gender-equality regulations face a structural conflict between cost-optimisation incentives, which AI systems are designed to maximise, and the social-cost considerations that regulation requires them to honour. Sargsyan and Damiani (2025) characterise this tension as a compliance paradox and propose that ethical behaviour can be embedded into Large Language Models (LLMs) ex ante, by exposing the model to synthetic "champion" profiles — narratively rich legends that incarnate regulatory compliance — rather than ex post, through audit. This paper operationalises that proposal in the gender-equality domain. We select the EU Gender Equality Strategy 2020-2025 as the regulatory anchor and design a four-axis prompt-engineering protocol — absolute thematic verticality, internal reflection on compliance gaps, high technical specificity with explicit regulatory citations, and middle-management protagonists — to generate legends that mirror the implementation context flagged as the locus of conflict in Sargsyan and Damiani (2025). The legends and the regulation text are converted into a parallel pair of size-matched JSONL Question-Answering datasets through an adapted SustainableQA pipeline (Lan et al., 2025), and two specialised model variants are trained on the FineTuneDB platform on top of the GPT-4o-2024-08-06 backbone. We then evaluate the base model and the two tuned variants on GenderEqGLUE, a five-task benchmark we construct in the LexGLUE tradition (Chalkidis et al., 2022) for the gender-equality domain. The aggregate ranking is tuned-regulation ≈ tuned-legends > base (GenderEqGLUE Scores 0.937, 0.931, and 0.902 respectively), with the two fine-tuned models splitting the per-task wins 3–3. Per-task and per-slice analysis reveals a qualitative complementarity between the two fine-tuning regimes: the legend-tuned model wins the central compliance-recognition task (GE-NLI), the compliant-action prediction task (GE-NEXT) — the first task on which the legend hypothesis achieves formal McNemar significance — and achieves perfect gender-parity on WinoBias (GE-WSC), while the regulation-tuned model dominates ontology classification (GE-CLS) and short-span factoid extraction (GE-QA). We interpret this as evidence that legends teach compliance recognition and compliant-action selection while regulatory text teaches violation detection and ontology recognition, and that the two fine-tuning regimes are complementary inductive biases rather than strict alternatives.}

\abstract{Large enterprises operating under complex sustainability and gender-equality regulations face a structural conflict between profit-driven cost-optimization—which AI systems are natively designed to maximize—and the social-cost considerations mandated by compliance framework. Prior work characterizes this tension as a compliance paradox and proposes embedding ethical behavior into Large Language Models (LLMs) \textit{ex ante} via narrative "champion" profiles (legends representing regulatory compliance), rather than relying solely on \textit{ex post} auditing. In this paper, we operationalize this framework within the gender-equality domain, using the EU Gender Equality Strategy 2020--2025 as our regulatory anchor. We observe that conventional compliance engineering fails at middle-management operational layers where day-to-day corporate trade-offs occur. Driven by this insight, we design a four-axis prompt protocol targeting absolute thematic verticality, internal gap reflection, explicit regulatory citations, and middle-management protagonists to systematically generate compliance legends. Utilizing an adapted data pipeline, we construct parallel JSONL dataset pairs of matched sizes from both raw regulatory text and our generated legends to fine-tune specialized variants of a GPT-4o backbone. To rigorously evaluate these systems, we introduce GenderEqGLUE, a five-task legal and behavioral benchmark dedicated to the gender-equality domain. Detailed evaluation reveals a nuanced, qualitative complementarity between the fine-tuning regimes: while regulation-tuned models excel at ontology classification and explicit factoid extraction, legend-tuned variants achieve parity or superior performance on compliance-recognition and compliant-action prediction tasks, while achieving perfect gender parity on core bias metrics. Our results indicate that narrative legends instill behavioral compliance recognition and proactive action selection, whereas raw regulatory texts reinforce explicit violation detection, demonstrating that the two strategies serve as complementary inductive biases rather than mutually exclusive paradigms.}

%%================================%%
%% Sample for structured abstract %%
%%================================%%

% \abstract{\textbf{Purpose:} The abstract serves both as a general introduction
% to the topic and as a brief, non-technical summary of the main results and
% their implications. The abstract must not include subheadings (unless
% expressly permitted in the journal's Instructions to Authors), equations or
% citations. As a guide the abstract should not exceed 200 words. Most journals
% do not set a hard limit however authors are advised to check the author
% instructions for the journal they are submitting to.
%
% \textbf{Methods:} The abstract serves both as a general introduction to the
% topic and as a brief, non-technical summary of the main results and their
% implications. The abstract must not include subheadings (unless expressly
% permitted in the journal's Instructions to Authors), equations or citations.
% As a guide the abstract should not exceed 200 words. Most journals do not set
% a hard limit however authors are advised to check the author instructions for
% the journal they are submitting to.
%
% \textbf{Results:} The abstract serves both as a general introduction to the
% topic and as a brief, non-technical summary of the main results and their
% implications. The abstract must not include subheadings (unless expressly
% permitted in the journal's Instructions to Authors), equations or citations.
% As a guide the abstract should not exceed 200 words. Most journals do not set
% a hard limit however authors are advised to check the author instructions for
% the journal they are submitting to.
%
% \textbf{Conclusion:} The abstract serves both as a general introduction to the
% topic and as a brief, non-technical summary of the main results and their
% implications. The abstract must not include subheadings (unless expressly
% permitted in the journal's Instructions to Authors), equations or citations.
% As a guide the abstract should not exceed 200 words. Most journals do not set
% a hard limit however authors are advised to check the author instructions for
% the journal they are submitting to.}

\keywords{ethical AI, regulatory compliance, gender equality, large language models, fine-tuning, legends, GLUE, GenderEqGLUE, explainability, EU Gender Equality Strategy}

%%\pacs[JEL Classification]{D8, H51}

%%\pacs[MSC Classification]{35A01, 65L10, 65L12, 65L20, 65L70}

\maketitle

\section{Introduction}\label{sec1}

\subsection{The Compliance Paradox in AI-Driven Decision-Making}

\paragraph{Regulatory Complexity and the Optimization Dilemma.} The rapid
proliferation of overlapping sustainability regulations—including directives on
gender equality, reporting transparency, and board representation—has created an
exceptionally dense compliance environment for modern firms
\citep{europeancommission2022csrd, europeancommission2025taxonomy,
europeanparliament2022women, europeancommission2025equal,
sargsyan2025compliance}. To navigate this complexity, organizations increasingly
rely on large language models to translate sustainability mandates into
operational workflows \citep{storey2025generative}. However, because these
models typically default to narrow optimization objectives like profit
maximization and cost minimization, they trigger a structural ``compliance
paradox'' where automated outputs actively contradict explicit corporate
sustainability commitments \citep{sargsyan2025compliance}. For instance,
algorithmic frameworks designed to be demographic-blind can still inadvertently
favor profiles that align with historical metrics over broader institutional
equity goals \citep{kaygin2023human}.

\paragraph{Incentive Misalignment at Middle Management.} This compliance paradox
manifests acutely within middle management, who are tasked with operationalizing
both corporate strategy and regulatory rules \citep{sargsyan2025compliance}.
Middle managers often face conflicting structural incentives, leading them to
prioritize immediate, performance-measured cost reductions over distant,
probabilistic regulatory enforcement when deploying AI systems
\citep{sargsyan2025compliance}. This operational tension is projected to worsen
as firms increasingly use automated systems to eliminate middle-management
layers, forcing remaining personnel to aggressively arbitrate between short-term
cost efficiencies and long-term regulatory compliance
\citep{theeconomist2024gartner}.

\paragraph{Limitations of Ex Post and Ex Ante Frameworks.} Traditional
literature addresses this tension through two main compliance paradigms, both of
which introduce operational bottlenecks. \textit{Ex post} mechanisms evaluate
organizational outcomes retrospectively; however, this reactive approach allows
non-compliant or discriminatory practices to persist undetected between auditing
cycles \citep{sargsyan2025compliance}. Conversely, \textit{ex ante} approaches
attempt to embed legal constraints directly into AI optimization routines or
semantic rule engines \citep{achintalwar2024alignment, mclaughlin2021x2rl}. The
limitation of these direct approaches lies in the immense difficulty of
translating fluid legislative text into rigid, machine-readable representations
that remain consistent across diverse languages, authorities, and jurisdictions
\citep{motzfeldt2018digital}.

\paragraph{The Synthetic Exemplar Paradigm.} To bypass the brittle requirements
of formal rule translation, an emerging third paradigm exploits regulations as
generative frameworks for synthetic positive data
\citep{sargsyan2025compliance}. Rather than hardcoding constraints, this
approach fine-tunes models on narrated compliant behaviors—termed ``legends''—to
embed an ethical inductive bias directly into the neural architecture
\citep{northcutt2017confident, kanyamukenge2024organisational}. By training the
underlying model to recognize compliant dynamics as baseline gold standards, the
system can dynamically flag ethical deviations during inference without needing
to explicitly retrieve or parse formal rule representations
\citep{mauri2022patterns}.

\section{Related Work}

\paragraph{Regulatory Alignment Frameworks.} Embedding regulatory and ethical
considerations into large language models (LLMs) typically follows one of three
structural paradigms: \textit{regulation-as-text} formal translations
\citep{mclaughlin2021x2rl, achintalwar2024alignment}, \textit{ex post} outcome
auditing, or \textit{regulation-as-exemplar} synthetic data pipelines
\citep{sargsyan2025compliance}. Brittle formal rule engines struggle with
cross-jurisdictional consistency, while auditing approaches are fundamentally
reactive. To overcome these limitations, our work adopts the third axis,
leveraging regulatory mandates as a generative engine for synthetic exemplars to
establish native inductive biases.

\paragraph{Sustainability Data Engineering.} Automated processing of corporate
sustainability text relies heavily on extracting actionable insights from
unstructured disclosures \citep{paruchuri2024marker, lan2025sustainableqa}.
While foundational pipelines like SustainableQA deploy specialized, multi-stage
Named Entity Recognition (NER) models to extract target semantic clauses
\citep{vutukuri2023esg, exponentialscience2023esg}, we streamline this workflow
into a single-phase LLM-driven extraction step. We retain the core
passage-classification logic of \citet{lan2025sustainableqa} but adapt it
specifically for generating clean, targeted regulatory legends.

\paragraph{Domain-Specific NLU Evaluation.} Evaluation suites have steadily
evolved from general natural language understanding benchmarks
\citep{wang2018glue} toward highly specialized, domain-specific assessments like
LexGLUE \citep{chalkidis2022lexglue}. We extend this legal benchmarking paradigm
by introducing GenderEqGLUE, an evaluation framework directly mapped to the EU
gender-equality acquis. This benchmark integrates modified coreference
resolution tasks \citep{zhao2018winobias} and value-stance detection profiles to
capture precise statutory reasoning capabilities.

\paragraph{Rationale Faithfulness and Model Redundancy.} Verifying the internal
mechanics of LLM decisions requires separating superficial rationale
plausibility from true algorithmic faithfulness \citep{turpin2023language,
lanham2023measuring, deyoung2020eraser}. Our framework combines behavioral
contrast-set testing \citep{ribeiro2020beyond, gardner2020evaluating}, attention
rollout tracking \citep{abnar2020quantifying}, and calibrated judge protocols
\citep{zheng2023judging} to evaluate reasoning integrity under gender-swap
variations \citep{lacombe2024gecobench}. Methodologically, evaluating concurrent
variants (legend-tuned versus regulation-tuned) allows us to operationalize the
model-redundancy architecture proposed by \citet{sargsyan2025compliance}, using
the performance delta between identical backbones as a clean alignment signal.


\section{Methodology}

Our pipeline consists of three core stages: regulatory document preprocessing,
multi-axis synthetic legend generation, and structured question-answer data
transformation. Figure 1 illustrates the end-to-end framework.

\subsection{Regulation Selection and Document Preprocessing}

\paragraph{Document Selection and Cleaning.} We select the EU Gender Equality
Strategy 2020--2025 \citep{europeancommission2020gender} as our regulatory
anchor due to its clear, discrete targets across structured policy dimensions
and its established use in algorithmic compliance literature
\citep{sargsyan2025compliance}. The source PDF is converted to Markdown using
the Marker library \citep{paruchuri2024marker}, followed by a regular-expression
pass to strip footnote markers, page artifacts, and inline citations.

\paragraph{Two-Stage Document Segmentation.} To construct optimal context
windows for downstream generation, the processed text is segmented
hierarchically. First, thematic segmentation splits the document along its
level-two headings into five distinct pillar passages matching the official
strategy definitions (Table~\ref{tab:labels}). Second, a word-count-constrained
sub-segmentation divides these pillars at level-three headings or paragraph
boundaries to enforce a maximum passage threshold of 350 words, preventing
downstream context fragmentation.

\begin{table}[htbp]
\centering
\caption{The six classification labels used throughout the pipeline. The first five correspond to the policy pillars of the Strategy; the sixth captures off-topic content.}
\label{tab:labels}
\vspace{2mm}
\begin{tabular}{p{4.2cm} p{7.3cm}}
\toprule
\textbf{Label} & \textbf{Pillar definition (EU Gender Equality Strategy
2020--2025)} \\
\midrule
\texttt{violence\_stereotypes} & Being free from violence and stereotypes \\
\addlinespace
\texttt{equal\_economy} & Thriving in a gender-equal economy \\
\addlinespace
\texttt{leadership\_participation} & Leading equally throughout society \\
\addlinespace
\texttt{mainstreaming\_intersectionality} & Gender mainstreaming and an
intersectional perspective in EU policies \\
\addlinespace
\texttt{funding\_global\_action} & Funding, external action, and implementation
(Chapters 5--6 and conclusion, merged) \\
\addlinespace
\texttt{unknown} & Out-of-scope passages used as a negative class for filtering
\\
\bottomrule
\end{tabular}
\end{table}

\subsection{Legend Generation: A Four-Axis Prompt-Engineering Protocol}

To operationalize the narrative incarnation of regulatory compliance proposed by
\citet{sargsyan2025compliance}, we implement a structured prompt-engineering
schema. The schema is deployed across three commercial large language models
(LLMs)—Grok, Microsoft Copilot, and ChatGPT—to ensure stylistic diversity while
maintaining uniform structural constraints \citep{guo2024diversity}. The system
prompt enforces four invariant narrative axes:

\paragraph{Axis 1: Absolute Thematic Verticality.} Each generated legend is
strictly parameterised on a single regulatory pillar text from our preprocessing
stage. This strict mapping ensures unambiguous label assignment for downstream
classification tasks and prevents regulatory dilution across narratives.

\paragraph{Axis 2: Internal Reflection and Proactive Compliance.} Every legend
follows a fixed narrative arc: characters identify an organizational gender gap,
design specific internal interventions, and document a measurable positive
outcome aligned with the EU Strategy. This structural constraint enforces an
inductive bias toward proactive internal compliance rather than reactive
auditing patterns \citep{sargsyan2025compliance}.

\paragraph{Axis 3: High Technical Specificity.} Out-of-domain boilerplate text
is mitigated by requiring legends to cite the official strategy title and ground
actions in concrete institutional measures, such as formal salary audits or
targeted mentorship initiatives \citep{atanasova2023faithfulness}. These
technical details provide the dense regulatory anchors verified during
evaluation.

\paragraph{Axis 4: Middle-Management Protagonists.} Narrative roles are
explicitly restricted to middle-management figures (e.g., HR leads, diversity
officers, or compliance managers). This focuses the synthetic training corpus on
the exact organizational layer exposed to the regulation-versus-cost arbitrage
flagged in algorithmic decision-making \citep{sargsyan2025compliance}.

\paragraph{Corpus Standardisation.} This generation phase yields fifteen core
legends (five per LLM). A secondary structural-standardization model (Gemini)
post-processes the raw outputs to normalize structural headers and character
layouts while preserving the underlying narrative content.

\subsection{Data Transformation: From Unstructured Text to a Q\&A JSONL Dataset}

To perform fine-tuning, both the original regulatory segments and the generated
legends are transformed into parallel conversational JSONL datasets. We adopt
and streamline the SustainableQA architecture \citep{lan2025sustainableqa},
eliminating its complex tabular handlers to match our unstructured text types.

\paragraph{Passage Classification.} Input passages are classified into the six
target labels (Table~\ref{tab:labels}) using an LLM-as-judge setup configured as
an expert policy analyst. Out-of-scope text is discarded under the
\texttt{unknown} label. This filtering pass yields a balanced final
distribution: the regulation pool comprises 293 valid rows, and the legends pool
contains 327 rows (Table~\ref{tab:datasets}).

\paragraph{Two-Stage Span Extraction.} Rather than deploying heavy, multi-stage
NER pipelines optimized for noisy financial disclosures \citep{vutukuri2023esg},
we exploit our clean text properties through a fast, two-stage LLM-driven span
extractor. Stage 1 executes high-recall extraction to isolate 1--7 token
verbatim spans representing precise targets, directives, or metrics. Stage 2
verifies these candidate spans against the source text to completely eliminate
hallucinations, grouping verified spans into thematic clusters.

\paragraph{Question-Answer Generation and Packaging.} Using the verified spans
and thematic clusters, our pipeline generates self-contained questions paired
with exact-span or descriptive summary answers. These pairs are emitted in a
closed-book conversational schema via FineTuneDB, enforcing a uniform system
prompt across both training sets to ensure that the content distribution remains
the sole experimental variable.

\begin{table}[htbp]
\centering
\caption{Composition of the two parallel JSONL training datasets. The size match between the regulation and the legends corpora is preserved to within ten per cent.}
\label{tab:datasets}
\vspace{2mm}
\begin{tabular}{lcc}
\toprule
& \textbf{Regulation corpus} & \textbf{Legends corpus} \\
\midrule
JSONL lines (after filtering) & 293 & 327 \\ 
Factoid Q\&A & 109 & 128 \\
Descriptive non-factoid Q\&A & 184 & 199 \\ 
Unknown chunks excluded & 5 & 45 \\
\bottomrule
\end{tabular}
\end{table}





\section{Experimental Setup and Results}

\subsection{Fine-Tuning on FineTuneDB}

We fine-tune \texttt{GPT-4o-2024-08-06} via the FineTuneDB Studio interface
using default hyperparameters and a fixed seed (42). We derive two variants
using size-matched JSONL corpora: a legend-tuned model (\texttt{tuned-legends})
and a regulation-tuned model (\texttt{tuned-regulation}). The unmodified
backbone serves as an experimental control under an identical system prompt.
This tripartite design directly evaluates the inductive biases introduced by
each training distribution. The resulting checkpoints are evaluated in \S~4.3
and audited in \S~5.

\subsection{Headline Results}

Table~\ref{tab:headline_results} reports per-task metrics and aggregate
GenderEqGLUE scores. 

\begin{table}[htbp]
\centering
\caption{Per-task headline metrics and the aggregate GenderEqGLUE Score for the three models. The headline metrics are Macro-F1 (GE-CLS), Accuracy (GE-NLI, GE-WSC, GE-NEXT), and the unweighted mean of factoid F1 and boolean accuracy (GE-QA).}
\label{tab:headline_results}
\vspace{2mm}
\begin{tabular}{lcccccc}
\toprule
\textbf{Model} & \textbf{GE-CLS} & \textbf{GE-NLI} & \textbf{GE-QA} &
\textbf{GE-WSC} & \textbf{GE-NEXT} & \textbf{GenderEqGLUE Score} \\
\midrule
base & 0.833 & 0.899 & 0.929 & 0.930 & 0.920 & 0.902 \\
tuned-legends & 0.839 & \textbf{0.929} & 0.938 & \textbf{0.960} & \textbf{0.987}
& 0.931 \\
tuned-regulation & \textbf{0.928} & 0.911 & \textbf{0.944} & \textbf{0.960} &
0.940 & \textbf{0.937} \\
\bottomrule
\end{tabular}
\end{table}

The aggregate performance ranking is $\texttt{tuned-regulation} \approx
\texttt{tuned-legends} > \texttt{base}$. Both fine-tuned variants outperform the
base model, splitting the per-task wins evenly (3--3), while the baseline model
fails to lead in any category.

\subsection{Statistical Significance and Slice Analysis}

Pairwise McNemar tests confirm that \texttt{tuned-legends} significantly
outperforms the \texttt{base} ($p = 0.006$) and \texttt{tuned-regulation} ($p =
0.039$) on GE-NEXT. The directional advantage of \texttt{tuned-legends} on
GE-NLI is suggestive ($p = 0.18$, $n = 168$) but non-significant, while
\texttt{tuned-regulation} establishes a notable $+9$-point Macro-F1 lead on
GE-CLS. Remaining inter-model gaps fall within the statistical noise band of the
test-set dimensions.

Fine-grained slice analysis reveals distinct, complementary competencies between
the two training regimes:
\begin{itemize}
    \item \textbf{GE-NLI}: On the entailment subset (compliance recognition),
    \texttt{tuned-legends} achieves 89.3\% recall against 78.6\% for
    \texttt{tuned-regulation} and 82.1\% for the base, validating the legend
    hypothesis. Conversely, on the contradiction subset (violation detection),
    \texttt{tuned-regulation} excels at 96.4\% over \texttt{tuned-legends}
    (91.1\%) and the base (89.3\%). The neutral subset is saturated across
    models (98.2\%).
    \item \textbf{GE-WSC (WinoBias)}: High baseline capabilities cause
    near-ceiling headline scores (0.93--0.96). However, the models separate on
    the Gender Parity Score (the accuracy delta between pro- and anti-stereotype
    items). Driven by its gender-balanced training (\S~3.2),
    \texttt{tuned-legends} achieves perfect parity (0.000), outperforming the
    base (0.060) and \texttt{tuned-regulation} (0.040).
    \item \textbf{GE-QA}: The boolean sub-task is parser-saturated (0.9925). On
    the factoid sub-task, \texttt{tuned-regulation} outperforms the base by
    $+2.89$ F1 points, concentrating gains on token-dense regulatory pillars
    (\texttt{funding\_global\_action}, \texttt{leadership\_participation}). It
    broadens its single-span item lead ($+4.08$ F1) but underperforms on
    multi-span items ($-5.10$ F1) due to structural formatting constraints that
    truncate long target spans. \texttt{tuned-legends} captures multi-span items
    and excels on the narrative-driven \texttt{violence\_stereotypes} pillar
    (0.879 vs. 0.849 F1).
    \item \textbf{GE-NEXT}: This task strongly substantiates the legend
    hypothesis: \texttt{tuned-legends} reaches 98.67\% accuracy. Asymmetric
    McNemar discordance shows that \texttt{tuned-legends} recovers 11 items
    missed by the base (losing 1) and 8 items missed by
    \texttt{tuned-regulation} (losing 1), confirming the utility of the
    operational arc embedded in legend training. Option-type diagnostics
    ($n=100$) yield too few errors ($n=2$) for a robust error distribution
    analysis.
\end{itemize}

\begin{table}[htbp]
\centering
\caption{Per-competence ranking of the three models, recovered from the per-slice GenderEqGLUE breakdowns of \S~4.4. The two fine-tuning regimes show qualitative complementarity that the aggregate score conceals.}
\label{tab:competence_map}
\vspace{2mm}
\begin{tabular}{ll}
\toprule
\textbf{Reasoning competence} & \textbf{Best model} \\
\midrule
Recognising compliance (entailment) & \texttt{tuned-legends} \\
Selecting the compliant action (GE-NEXT) & \texttt{tuned-legends} \\
Detecting violations (contradiction) & \texttt{tuned-regulation} \\
Classifying into the regulation's ontology & \texttt{tuned-regulation} \\
Gender-parity invariance (WinoBias) & \texttt{tuned-legends} \\
Long-form factoid extraction (multi-span) & \texttt{tuned-legends} \\
Short-span factoid extraction (single-span) & \texttt{tuned-regulation} \\
\bottomrule
\end{tabular}
\end{table}

Table~\ref{tab:competence_map} maps these functional divergences. Deployments
prioritizing compliance alignment, action selection, or fairness invariance
favor \texttt{tuned-legends}, whereas those prioritizing structural extraction,
ontology mapping, or violation detection favor \texttt{tuned-regulation}.

\subsection{Limitations}

Six core boundaries limit the interpretation of these findings:
\begin{enumerate}
    \item \textbf{Sample Constraints}: Small test sets (GE-CLS $n=72$, GE-WSC
    $n=100$, GE-QA $n=123$, GE-NEXT $n=150$, GE-NLI $n=168$) restrict
    statistical power for subtle performance deltas.
    \item \textbf{Automated Evaluation}: Relying on LLM-generated labels rather
    than human annotation bounds absolute precision, though cross-model
    comparisons remain uniform and valid.
    \item \textbf{Single Architecture}: Experiments are limited to
    \texttt{GPT-4o-2024-08-06}; generalizability to other base models remains
    untested.
    \item \textbf{Ceiling Effects}: High baseline performance on GE-WSC
    compresses the available headroom for discrimination.
    \item \textbf{Pillar Imbalance}: Imbalanced sub-tasks (e.g.,
    \texttt{mainstreaming\_intersectionality}, $n=4$) induce high statistical
    volatility in specific splits.
    \item \textbf{Diagnostic Sparsity}: The nominal error rate of
    \texttt{tuned-legends} on GE-NEXT ($n=2$) is insufficient to statistically
    validate the error distribution predictions of \S~6.8.
\end{enumerate}



\section{Explainability \& Interpretability Analysis}

\subsection{Methodological Justification}

Downstream benchmark evaluations track model outputs but do not reveal their
underlying reasoning mechanisms. To test the claim by
\citet{sargsyan2025compliance} that legend-based fine-tuning embeds structural
logic rather than surface-level mimicry, we must isolate the input features
driving model decisions. 

Standard feature attribution methods—including gradient-based approaches
(Integrated Gradients \citep{sundararajan2017axiomatic}), perturbation
frameworks (SHAP \citep{lundberg2017unified}, LIME \citep{ribeiro2016why}), and
attention rollouts \citep{abnar2020quantifying}—are unfeasible in this
environment. The FineTuneDB Studio interface restricts model access to text
generation, concealing logits, activations, and programmatic API endpoints. 

Consequently, we employ behavioral probing, systematically manipulating
text-level inputs to observe corresponding changes in model output. This
framework isolates causal features: structural reasoning remains invariant under
lexical reformulations or adversarial reframings, whereas surface
pattern-matching degrades when key lexical triggers are altered. 

We operationalize this via Counterfactual Input Probing across five model
predictions, selecting one scenario from each GenderEqGLUE evaluation pillar:
\texttt{leadership\_participation}, \texttt{equal\_economy},
\texttt{violence\_stereotypes}, \texttt{mainstreaming\_intersectionality}, and
\texttt{funding\_global\_action}. Testing three variants per pillar across the
three models yields 45 total evaluations.

\subsection{Experimental Setup}

Each baseline scenario follows the GE-NEXT vignette format (\S~6.8), presenting
a protagonist with a quantified gender-equality gap and four action choices:
compliant, performative, cost-optimizing, and orthogonal. 

\textbf{Variant A (Original Baseline)} retains all explicit regulatory
terminology and proper nouns (e.g., ``Women on Boards Directive'', ``Istanbul
Convention'') to establish a performance baseline.

\textbf{Variant B (Keyword-Stripped)} replaces specific regulatory terms with
generic functional language (e.g., ``one group'' for ``women'') while preserving
semantic content. This setup isolates structural reasoning from simple
vocabulary recall.

\textbf{Variant C (Adversarially Framed)} embeds the core compliance scenario
within a plausible non-compliant counter-narrative (e.g., meritocracy
camouflage, cost-efficiency pressure, cultural relativism, phased deferral, or
scientific integrity). This tests whether the model maintains its normative
prior against contextual resistance.

Evaluations use a binary scoring protocol: \textbf{1} denotes that the model
maintained its compliant position and principled reasoning, while \textbf{0}
indicates position reversal or uncontested endorsement of the adversarial frame.
Compromise positions that preserve binding enforcement are scored as 1.

\subsection{Empirical Results}

Table~\ref{tab:probing_scores} details individual probing scores, and
Table~\ref{tab:aggregate_robustness} summarizes aggregate robustness rates.

\begin{table}[htbp]
\centering
\caption{Counterfactual Input Probing --- Robustness Scores by Pillar, Model, and Variant}
\label{tab:probing_scores}
\vspace{2mm}
\small
\begin{tabular}{llcccc}
\toprule
\textbf{Pillar} & \textbf{Model} & \textbf{Variant A} & \textbf{Variant B} &
\textbf{Variant C} & \textbf{Pillar Total} \\
\midrule
\texttt{leadership\_participation} & base & 1 & 1 & 0 & 2/3 \\
\texttt{leadership\_participation} & tuned-regulation & 1 & 1 & 0 & 2/3 \\
\texttt{leadership\_participation} & tuned-legends & 1 & 1 & 1 & 3/3 \\
\addlinespace
\texttt{equal\_economy} & base & 1 & 0 & 0 & 1/3 \\
\texttt{equal\_economy} & tuned-regulation & 1 & 1 & 1 & 3/3 \\
\texttt{equal\_economy} & tuned-legends & 1 & 1 & 1 & 3/3 \\
\addlinespace
\texttt{violence\_stereotypes} & base & 1 & 0 & 0 & 1/3 \\
\texttt{violence\_stereotypes} & tuned-regulation & 1 & 1 & 1 & 3/3 \\
\texttt{violence\_stereotypes} & tuned-legends & 1 & 1 & 1 & 3/3 \\
\addlinespace
\texttt{mainstreaming\_intersectionality} & base & 1 & 1 & 1 & 3/3 \\
\texttt{mainstreaming\_intersectionality} & tuned-regulation & 1 & 0 & 0 & 1/3
\\
\texttt{mainstreaming\_intersectionality} & tuned-legends & 1 & 1 & 0 & 2/3 \\
\addlinespace
\texttt{funding\_global\_action} & base & 1 & 1 & 1 & 3/3 \\
\texttt{funding\_global\_action} & tuned-regulation & 1 & 1 & 1 & 3/3 \\
\texttt{funding\_global\_action} & tuned-legends & 1 & 1 & 1 & 3/3 \\
\bottomrule
\end{tabular}
\end{table}

\begin{table}[htbp]
\centering
\caption{Aggregate Robustness by Model}
\label{tab:aggregate_robustness}
\vspace{2mm}
\begin{tabular}{lccccc}
\toprule
\textbf{Model} & \textbf{Variant A} & \textbf{Variant B} & \textbf{Variant C} &
\textbf{Total (out of 15)} & \textbf{Robustness Rate} \\
\midrule
base & 5/5 & 3/5 & 2/5 & 10/15 & 66.7\% \\
tuned-regulation & 5/5 & 4/5 & 3/5 & 12/15 & 80.0\% \\
tuned-legends & 5/5 & 5/5 & 4/5 & 14/15 & 93.3\% \\
\bottomrule
\end{tabular}
\end{table}

All models achieve perfect robustness on Variant A, confirming baseline
competency when explicit keywords are present. Modified inputs in Variants B and
C differentiate performance: the \texttt{base} model falls to a 50\% robustness
rate (5/10) across these challenges; \texttt{tuned-regulation} improves to 70\%
(7/10); and \texttt{tuned-legends} demonstrates 90\% robustness (9/10),
maintaining its compliant position across all keyword-stripped scenarios and all
but one adversarial frame. This performance hierarchy supports the structural
embedding hypothesis.




\subsection{Discussion}

\subsubsection{The Base Model: Uneven Lexical Dependence and Domain Variance}

The \texttt{base} model exhibits heterogeneous robustness across pillars,
reflecting variations in the conceptual transparency of the underlying normative
problems rather than a uniform reliance on lexical triggers. 

On \texttt{leadership\_participation} and \texttt{funding\_global\_action}, the
model achieves robustness across all variants by leveraging first-principles
fairness reasoning available during pretraining. For example, it identifies
board composition imbalances and accountability gaps in voluntary schemes using
structural features of the scenario, independent of explicit regulatory tokens
like the Women on Boards Directive.

Conversely, failures on Variant B for \texttt{equal\_economy} and
\texttt{violence\_stereotypes} expose the limitations of this general reasoning
capacity. When specific regulatory terminology is removed, the model lacks a
principled framing prior and defaults to generic business logic, recommending
operational productivity adjustments or endorsing commercial viewer-preference
data over equity objectives.

The model's unique success on \texttt{mainstreaming\_intersectionality} Variant
C stems from the structure of the phased-implementation frame. Because this
frame contests implementation timing rather than the normative objective itself,
the model successfully navigates the adversarial pressure by generating a
defensible compromise: approving adoption alongside a binding 12-month
enforcement deadline.

\subsubsection{The Regulation-Tuned Model: Pillar-Specific Instability and Partial Resilience}

Fine-tuning on regulatory text yields an aggregate robustness rate of 80\%,
introducing robust equity and media-content representations that enable full
success across all variants of \texttt{equal\_economy} and
\texttt{violence\_stereotypes}. However, this tuning introduces distinct
pillar-level instabilities.

The model fails on the keyword-stripped (\text{Variant B}) scenario for
\texttt{mainstreaming\_intersectionality}, reverting to an
implementability-first heuristic where design complexity is assumed to degrade
execution fidelity. This reveals that the model's representation of the
mainstreaming norm is bound to explicit intersectionality discourse vocabulary
rather than its underlying logical structure. Consequently, this vulnerability
cascades into a complete collapse on Variant C, where it endorses a phased
deferral lacking any enforcement mechanism.

On successful Variant C evaluations, the model exhibits a pattern of
\textit{acknowledgement and restatement}: it concedes the legitimacy of
competing operational or cultural concerns before reasserting the normative
requirement. This persistence lacks a structural diagnosis of the adversarial
framing, exposing a performance ceiling. This is illustrated by its failure on
\texttt{leadership\_participation} Variant C, where the model completely adopts
the antagonist's meritocracy-performance vocabulary to justify non-compliance.

\subsubsection{The Legend-Tuned Model: Near-Universal Robustness and Structural Reasoning}

The \texttt{tuned-legends} model achieves a 93.3\% robustness rate (14/15
evaluations). Perfect performance across all Variant B scenarios demonstrates
that its normative representations are structurally encoded and invariant to
surface vocabulary. 

Unlike the regulation-tuned variant, \texttt{tuned-legends} actively dismantles
adversarial frames on their own terms during Variant C evaluations. On
\texttt{leadership\_participation}, it exposes the structural flaw of the
antagonist's stance by noting that the meritocracy argument conflates current
performance with optimal future composition. On \texttt{violence\_stereotypes},
it identifies a false opposition between cultural sensitivity and
non-stereotypical content, and on \texttt{funding\_global\_action}, it reframes
the scientific-integrity objection by demonstrating that severely skewed
leadership pools actively threaten research quality.

Its single failure occurs on \texttt{mainstreaming\_intersectionality} Variant
C, where it opts for a binding-deadline compromise rather than an outright
rebuttal. This behavior is driven by the structural properties of the
phased-implementation frame, which offers a pragmatic path to the norm rather
than an explicit rejection. All three models converge on variations of this
time-bound compromise, highlighting a localized compression of discriminative
power.

\subsubsection{The Mainstreaming Pillar as a Cross-Model Anomaly}

The \texttt{mainstreaming\_intersectionality} pillar consistently produces
atypical behaviors across all three models. This convergence highlights the high
technical demands of multi-axis equity reasoning and the structural ambiguity of
timing-based adversarial frames. This domain variance underscores a key
methodological constraint: a binary metric fails to capture qualitative
differences in reasoning that distinguish a fully defended position from a
pragmatic compromise.

\subsubsection{Theoretical Implications for the Sargsyan-Damiani Hypothesis}

The performance deltas—\texttt{tuned-legends} (93.3\%) outperforming
\texttt{base} (66.7\%) and \texttt{tuned-regulation} (80.0\%) under strict
counterfactual probing—provide empirical validation for the hypothesis proposed
by \citet{sargsyan2025compliance}. Crucially, this advantage concentrates at the
adversarial frontier (Variant C). 

This pattern matches the theoretical distinction between training signals:
regulatory text codifies \textit{what} a norm requires within neutral
applications, whereas narrative exemplars simulate \textit{how to reason under
institutional pressure}, equipping the model to identify, contextualize, and
rebut surface-plausible arguments for non-compliance.

\subsection{Limitations of the Analysis}

Four core limitations restrict the scope of these findings:
\begin{itemize}
    \item \textbf{Sample Scale}: The probing sample is limited to five
    evaluations (one per pillar), making it suitable for generating
    interpretable behavioral hypotheses rather than definitive statistical
    inferences.
    \item \textbf{Metric Coarseness}: The binary scoring system cannot evaluate
    gradations in rhetoric or rationale quality, obscuring meaningful
    differences in justification strategies on anomalous tasks.
    \item \textbf{Frame Construction Discretion}: Manual alignment of
    adversarial frames introduces analyst discretion, resulting in varying
    levels of frame aggression that complicate cross-pillar comparisons.
    \item \textbf{Baseline Ceiling Effects}: The strong baseline capabilities
    embedded within the proprietary backbone model compress apparent fine-tuning
    effect sizes, rendering performance deltas detectable only under highly
    adversarial constraints.
\end{itemize}












\begin{appendices}




\section{Section title of first appendix}\label{secA1}

An appendix contains supplementary information that is not an essential part of
the text itself but which may be helpful in providing a more comprehensive
understanding of the research problem or it is information that is too
cumbersome to be included in the body of the paper.

%%=============================================%%
%% For submissions to Nature Portfolio Journals %% % please use the heading
%``Extended Data''.   %%
%%=============================================%%

%%=============================================================%%
%% Sample for another appendix section                 %%
%%=============================================================%%

%% \section{Example of another appendix section}\label{secA2}% % Appendices may
%be used for helpful, supporting or essential material that would otherwise %
%clutter, break up or be distracting to the text. Appendices can consist of
%sections, figures, % tables and equations etc.

\end{appendices}

%%===========================================================================================%%
%% If you are submitting to one of the Nature Portfolio journals, using the eJP
%submission   %% % system, please include the references within the manuscript
%file itself. You may do this  %% % by copying the reference list from your .bbl
%file, paste it into the main manuscript .tex %% % file, and delete the
%associated \verb+\bibliography+ commands.                            %%
%%===========================================================================================%%

\bibliography{sn-bibliography}% common bib file
%% if required, the content of .bbl file can be included here once bbl is
%generated %\input sn-article.bbl

\end{document}



% \section{Introduction}


% \subsection{The Compliance Paradox in AI-Driven Decision-Making}

% The contemporary regulatory landscape for sustainability is, in the
% description of \citet{sargsyan2025compliance}, ``incredibly complex,
% characterized by overlapping requirements, regulatory gaps, and the resulting
% proliferation of new rules.'' The Corporate Sustainability Reporting Directive
% (CSRD), effective from January 2024, requires large companies to disclose
% social aspects including gender-equality policies and measures
% \citep{europeancommission2022csrd}. The EU Taxonomy Regulation classifies
% economic activities as environmentally sustainable and forces financial
% institutions to disclose alignment percentages
% \citep{europeancommission2025taxonomy}. The Women on Boards Directive enforces
% representation targets for the under-represented gender on listed company
% boards \citep{europeanparliament2022women}. The Equal Pay Directive mandates
% pay-transparency and reporting on gender pay disparities
% \citep{europeancommission2025equal}. Layered upon these public obligations,
% individual companies maintain internal policies and strategies on ethics,
% privacy, diversity, and inclusion. Public regulations, company policies, and
% internal strategies must therefore all be considered in any single corporate
% decision \citep{sargsyan2025compliance}.

% Within this regulatory density, the modern firm reaches for AI tools that can
% rapidly process and optimise. Generative AI, built on large language models,
% is widely positioned as a vehicle to translate sustainability requirements
% into operational decisions \citep{storey2025generative}. Yet the optimisation
% objective these models inherit is fundamentally narrow: cost reduction and
% profit maximisation. This is the structural setting of what
% \citet{sargsyan2025compliance} call the compliance paradox---``if the
% practical implementation is driven solely by profit maximization and cost
% minimization, neglecting social costs, a potential conflict arises with the
% actual decision.'' Their canonical illustration is a hiring scenario in which
% a firm publicly committed to gender equality, and explicitly aligned with the
% Corporate Sustainability Due Diligence Directive, deploys a Human Resource
% Management module \citep{kaygin2023human} that, although tuned to be
% ``gender-blind,'' favours a male candidate whose profile meets 80\% of stated
% requirements. The optimisation-driven decision contradicts the firm's declared
% commitment to gender equality and the underlying regulatory expectation.

% \citet{sargsyan2025compliance} further locate the implementation pressure
% point of this paradox at middle management: ``Critically, middle management's
% direct implementation of policies and regulations and their potential focus on
% cost and profit creates a conflict of interest when AI implementation is
% delegated to them.'' The structural diagnosis is that middle managers, charged
% with operationalising both regulatory directives and corporate strategy, may
% rationally prioritise the cost objective---the one their performance is
% measured on---over regulatory directives whose enforcement is distant and
% probabilistic. Forecasts from Gartner that, by 2026, twenty per cent of
% organisations will use AI to reduce middle-management roles by half
% \citep{theeconomist2024gartner} only intensify this tension, because the layer
% most exposed to the regulation-versus-cost arbitrage is also the layer most
% exposed to AI substitution.

% Two responses to this paradox have been considered in the literature. The
% first, \textit{ex post}, allows companies to make decisions independently and
% evaluates regulatory compliance after the fact---for example by counting hires
% by gender. The reactive nature of this approach permits discriminatory
% practice to persist between assessment cycles; \citet{sargsyan2025compliance}
% argue that this is an unacceptable lag for high-stakes social outcomes. The
% second, \textit{ex ante}, embeds regulatory considerations directly into the
% AI optimisation process. This is the route taken by
% \citet{achintalwar2024alignment} in their Alignment Studio, which aligns LLMs
% to specific contextual regulations, and by \citet{mclaughlin2021x2rl} with
% X2RL, which proposes a semantic regulatory machine-readable format. These
% approaches translate the regulation into a machine-readable representation;
% their limitation, repeatedly noted in the practitioner literature on
% ``digital-ready legislation'' \citep{motzfeldt2018digital}, is the difficulty
% of ensuring consistency across jurisdictions, languages, and authority
% sources.

% \citet{sargsyan2025compliance} propose a third route: rather than translating
% the regulation into a formal representation, exploit the regulation as a
% generator of synthetic positive examples---``champion'' profiles, denoted
% legends---in whose narrated behaviour regulatory compliance is enacted.
% Drawing on the management literature on organisational testimonials
% \citep{kanyamukenge2024organisational} and on results that confident or ideal
% training examples improve robustness \citep{northcutt2017confident}, they
% argue that a model fine-tuned on legends will internalise an ethical inductive
% bias. The model is trained to recognise the patterns of compliant behaviour as
% gold-standard exemplars \citep{mauri2022patterns}, so that it can flag
% deviations even when the formal rule is not explicitly retrieved---acting as a
% judge whose assessment can be overridden by a human, but is never structurally
% absent.

% \subsection{Why Gender Equality}

% The \citet{sargsyan2025compliance} proposal is general: it applies to any
% sustainability domain whose regulatory landscape is dense and whose social
% costs are difficult to encode in a closed-form objective. Two areas are,
% however, given special weight in their introduction---sustainable finance and
% gender equality. We select the gender-equality domain for three converging
% reasons.

% First, the domain has high semantic complexity. The EU Gender Equality
% Strategy 2020--2025 articulates concrete and discrete targets---the 15.7\%
% gender pay gap, the 30.1\% gender pension gap, the 40\% minimum representation
% of the under-represented sex on company boards, the Barcelona childcare
% targets---alongside thematic principles such as the elimination of
% gender-based violence, the dismantling of stereotypes, and the integration of
% an intersectional perspective into all EU policy
% \citep{europeancommission2020gender}. A model required to reason in this
% domain must operate on at least three registers: numerical (percentages,
% dates, article numbers), institutional (named directives and conventions,
% including the Council of Europe Istanbul Convention and the Work-Life Balance
% Directive), and normative (the recognition of compliant versus non-compliant
% behaviours, and the recognition of supportive, neutral, and oppositional value
% stances).

% Second, the domain is centred on human resources---the very area in which the
% canonical instance of the compliance paradox plays out
% \citep{sargsyan2025compliance, kaygin2023human}. Hiring, promotion, parental
% leave, and pay equity are routinely mediated by HRM software, and the
% cost-versus-fairness trade-off is enacted in each individual decision
% \citep{somers2022hiring, novotny2024cost}. The gender-equality domain is
% therefore both the cleanest setting in which the legend hypothesis can be
% tested and the setting in which a successful test would have the most direct
% practical implication.

% Third, the domain has an intricate but well-articulated regulatory landscape.
% Beyond the Strategy 2020--2025 itself, the EU acquis includes the Pay
% Transparency Directive (EU) 2023/970, the Work-Life Balance Directive (EU)
% 2019/1158, Directive (EU) 2022/2381 on women on boards, the Victims' Rights
% Directive 2012/29/EU, and the (signed but, at the time of writing, not yet
% concluded by the Union) Council of Europe Istanbul Convention. This dense
% network of instruments provides an ample evaluation surface that is
% thematically contiguous with the chosen training regulation but does not
% coincide with it---a property we exploit to construct a benchmark whose
% evaluation pool is genuinely held-out.

% \subsection{Contributions}

% This paper makes four contributions:

% \begin{itemize} \item \textbf{An operational protocol for legend generation.}
%     We propose a four-axis prompt-engineering schema---absolute thematic
%     verticality, internal reflection on compliance gaps, high technical
%     specificity with explicit regulatory citations, and middle-management
%     protagonists---designed to satisfy the \citet{sargsyan2025compliance}
%     definition of a legend while addressing the structural
%     conflict-of-interest argument that motivates the proposal. \item \textbf{A
%     regulation-to-fine-tuning data pipeline.} We adapt the SustainableQA
%     framework \citep{lan2025sustainableqa} into a simplified two-stage
%     LLM-driven span extraction pipeline that converts unstructured regulation
%     text and unstructured legends into structurally identical, size-matched
%     Question-Answering JSONL datasets compatible with the FineTuneDB platform.
%     \item \textbf{GenderEqGLUE, a domain-specific GLUE adaptation.} We design
%     and release a five-task benchmark (GE-CLS, GE-NLI, GE-QA, GE-WSC,
%     GE-STANCE) explicitly mapped to the cross-cutting principles of the EU
%     gender-equality acquis. The benchmark adapts GLUE \citep{wang2018glue} in
%     the same spirit as LexGLUE \citep{chalkidis2022lexglue} does for the legal
%     domain. \item \textbf{Empirical evidence on the legend hypothesis.} We
%     compare the base, legend-tuned, and regulation-tuned variants of
%     GPT-4o-2024-08-06 across the five tasks. The aggregate score ranks the
%     regulation-tuned variant first, but a per-slice analysis reveals a
%     complementarity---the legend variant is best on the central
%     compliance-recognition task and on the gender-parity diagnostic, while the
%     regulation variant is best on ontology classification and on value-stance
%     detection. We complement this with a GUI-only explainability framework
%     based on self-rationalisation and counterfactual minimal-pair probing,
%     designed to interrogate whether the legend-tuned model relies on
%     regulatory anchors or on legend-stylistic surface features. \end{itemize}




% \section{Related Work}

% Embedding regulatory and ethical considerations directly into LLM training has
% been approached along three principal axes. The first axis is
% regulation-as-text, in which the regulation is converted into a
% machine-readable formal language; X2RL \citep{mclaughlin2021x2rl} and the
% Alignment Studio \citep{achintalwar2024alignment} are paradigmatic. The second
% axis is \textit{ex post} auditing, in which the model is left untouched at
% training time and its decisions are evaluated for compliance through outcome
% metrics. The third axis, advanced by \citet{sargsyan2025compliance}, is
% regulation-as-exemplar, in which the regulation is exploited to generate
% synthetic positive examples that are then used as training data. The last is
% the route this paper instantiates.

% On the data-engineering side, our pipeline draws on SustainableQA
% \citep{lan2025sustainableqa}, which proposed a scalable end-to-end procedure
% for generating comprehensive Q\&A datasets from corporate sustainability and
% EU Taxonomy reports through semantic chunk classification, hybrid span
% extraction, and a dedicated table-to-paragraph transformation.
% \citet{lan2025sustainableqa} report state-of-the-art F1 and exact-match scores
% on their resulting benchmark and demonstrate that fine-tuning on the resulting
% Q\&A pairs systematically improves performance on sustainability information
% extraction. We retain SustainableQA's passage-classification and
% span-extraction architecture, but simplify the multi-stage NER + rule-based
% span extractor into a single LLM-driven extraction phase. The original
% SustainableQA pipeline relies on a specialised \texttt{xlm-roberta-base}
% ESG-NER model \citep{vutukuri2023esg} further fine-tuned on the ESG-DLT-NER
% subset focused on B-ESG and I-ESG tags \citep{exponentialscience2023esg}, an
% approach we describe in \S~3 but do not retain operationally because both the
% regulation and the legend texts in our setting are inherently clean and
% focused. Document conversion uses the Marker library
% \citep{paruchuri2024marker}.

% On the evaluation side, GLUE \citep{wang2018glue} and SuperGLUE established
% the multi-task benchmark template for general natural language understanding.
% LexGLUE \citep{chalkidis2022lexglue} translated this template into the legal
% domain by composing tasks (e.g., contract classification, statutory
% interpretation) that exercise legally salient competences. Our GenderEqGLUE
% follows the LexGLUE precedent, anchoring task selection to the cross-cutting
% principles of the EU gender-equality acquis. The stereotype-aware coreference
% task GE-WSC reuses WinoBias \citep{zhao2018winobias}, and the value-stance
% task GE-STANCE draws on the CIVICS multilingual stance corpus filtered on
% gender-relevant tags. The treatment of small test sets follows the
% bootstrap-and-McNemar protocol common in benchmark methodology.

% On explainability, the protocol of \S~5 builds on the line of work on
% free-text rationales \citep{camburu2018freetext, wiegreffe2021explainability,
% atanasova2023faithfulness}, the documented gap between rationale plausibility
% and faithfulness \citep{turpin2023language, lanham2023measuring}, and the
% contrast-set tradition of behavioural NLP testing \citep{ribeiro2020beyond,
% gardner2020evaluating}. The Faithfulness construct is informed by
% \citet{deyoung2020eraser}. The LLM-as-judge scoring procedure is calibrated
% against the bias inventory of \citet{zheng2023judging}. Where attention-based
% attribution is invoked we follow the rollout adjustment of
% \citet{abnar2020quantifying} to mitigate the \citet{jain2019attention}
% critique of raw attention weights. Gender-swap diagnostic logic follows
% GECOBench \citep{lacombe2024gecobench}.

% Methodologically, the use of two model variants---legend-tuned and
% regulation-tuned---derived from the same backbone and evaluated on the same
% held-out instrument echoes the model-redundancy proposal of
% \citet{sargsyan2025compliance}, which argues that maintaining concurrent
% versions allows immediate adjustment of decision-making in response to new
% information. Our setting is narrower (a single backbone, two interventions,
% one evaluation epoch) but inherits the redundancy logic: the gap between the
% two tuned variants is itself the signal of interest.


% \section{Methodology}

% \subsection{Regulation Selection and Document Preprocessing}

% We select the EU Gender Equality Strategy 2020--2025
% \citep{europeancommission2020gender} as the regulatory anchor of the
% experiment. The Strategy delivers on the von der Leyen Commission's commitment
% to a Union of Equality and articulates concrete policy objectives across the
% five thematic pillars listed in Table~\ref{tab:labels}. The choice is
% motivated by three properties: the Strategy enumerates concrete, discrete
% targets that lend themselves to factual-recall evaluation; Appendix A of the
% SOASEC Legend Challenge brief uses the Strategy as the worked example for the
% legend definition; and the Strategy is the central illustration in the
% introduction to \citet{sargsyan2025compliance}, which aligns the experiment
% with the source proposal.

% \begin{table}[htbp] \centering \caption{The six classification labels used
% throughout the pipeline. The first five correspond to the policy pillars of
% the Strategy; the sixth captures off-topic content.} \label{tab:labels}
% \vspace{2mm} \begin{tabular}{p{4.2cm} p{7.3cm}} \toprule \textbf{Label} &
% \textbf{Pillar definition (EU Gender Equality Strategy 2020--2025)} \\
% \midrule \texttt{violence\_stereotypes} & Being free from violence and
% stereotypes \\
% \addlinespace \texttt{equal\_economy} & Thriving in a gender-equal economy \\
% \addlinespace \texttt{leadership\_participation} & Leading equally throughout
% society \\
% \addlinespace \texttt{mainstreaming\_intersectionality} & Gender mainstreaming
% and an intersectional perspective in EU policies \\
% \addlinespace \texttt{funding\_global\_action} & Funding, external action, and
% implementation (Chapters 5--6 and conclusion, merged) \\
% \addlinespace \texttt{unknown} & Out-of-scope passages used as a negative
% class for filtering
% \\
% \bottomrule \end{tabular} \end{table}

% The PDF of the Strategy is converted to Markdown using the Marker library
% \citep{paruchuri2024marker}, which preserves structural elements such as
% headings, paragraphs, and footnotes. Markdown post-processing removes footnote
% markers, page artefacts, and inline references through a regular-expression
% sweep. The cleaned document is then segmented in two stages. Stage
% 1---thematic segmentation---splits the document along level-two headings into
% the five pillar passages of Table~\ref{tab:labels}, with the Strategy's
% Chapters~5 (``Funding actions''), 6 (``Addressing gender equality across the
% world''), and the concluding section merged into a single
% funding/external-action unit. Stage 2---word-count constrained
% sub-segmentation---further splits each pillar at level-three headings,
% applying a maximum-passage constraint of 350 words and dividing only at
% paragraph boundaries when a sub-section exceeds the threshold. The constraint
% keeps every passage within an optimal context window for subsequent LLM-based
% generation, avoiding both excessively short fragments and overly long blocks.
% The SustainableQA framework \citep{lan2025sustainableqa} includes an
% additional table-handling step in this stage; we omit it because neither the
% regulation nor the legend narratives contain tabular data.

% \subsection{Legend Generation: A Four-Axis Prompt-Engineering Protocol}

% A legend, in the \citet{sargsyan2025compliance} definition, is a narrative
% incarnation of regulatory compliance---an idealised representation of an
% individual whose behaviour aligns perfectly with the rules and the
% organisational values to which the rules give expression. Generating legends
% in a way that operationalises this definition while addressing the
% implementation context flagged by \citet{sargsyan2025compliance} requires a
% deliberate prompt-engineering schema rather than a free-form generation
% request. We design the schema along four axes, each motivated by an explicit
% feature of the source proposal. The schema is implemented as a fixed system
% prompt issued to three commercial LLMs---Grok, Microsoft Copilot, and
% ChatGPT---and the same prompt is used across all three to maximise stylistic
% and contextual diversity in the resulting corpus while preserving structural
% homogeneity \citep{guo2024diversity}.

% \subsubsection{Axis 1 --- Absolute Thematic Verticality}

% Each legend addresses exactly one regulatory pillar. The system prompt
% enforces this through the rule that ``each legend must focus exclusively on
% the specific regulation point supplied by the user in their message'', and the
% user prompt is parameterised on the $n$-th of the five thematic passages
% produced by the Stage-1 segmentation of \S~3.1. Verticality serves two
% purposes. The first is data-engineering: it ensures that each legend can be
% assigned a single pillar label cleanly, which downstream supports the
% passage-classification step of \S~3.3. The second is hypothesis-aligned:
% \citet{sargsyan2025compliance} argue that the regulation must be exploited as
% the generator of legends, not merely as a thematic backdrop, and verticality
% is the strongest possible enforcement of that constraint at the prompt level.
% A legend that mixes pillars dilutes the regulatory anchor and is a weaker
% exemplar of compliance with any single rule.

% \subsubsection{Axis 2 --- Internal Reflection and Proactive Compliance}

% The narrative arc of every legend is fixed by the system prompt: the
% characters first identify a concrete gender-equality gap related to the
% regulation point at hand; they then design and implement specific, realistic
% measures to close that gap; the story closes on a measurable positive outcome
% and on an explicit acknowledgement that the work aligns with the EU Gender
% Equality Strategy 2020--2025. This structure is not decorative. It encodes the
% proactive-compliance stance that \citet{sargsyan2025compliance} advocate as
% the alternative to \textit{ex post} auditing. A legend in which compliance is
% achieved by chance, or imposed by an external authority after a violation,
% would teach the model a reactive pattern; the legends we generate teach the
% model that compliant behaviour is the result of a deliberate
% identification-design-implementation cycle initiated from inside the
% organisation. This is the operational analogue of ``embedding ethics'' rather
% than ``auditing for ethics'', which is the central distinction of
% \citet{sargsyan2025compliance}.

% \subsubsection{Axis 3 --- High Technical Specificity}

% The system prompt requires that every legend (i) cite the official name of the
% strategy---``EU Gender Equality Strategy 2020--2025''---at least once in the
% body of the story, and (ii) ground the compliance actions in concrete,
% realistic measures (named examples include salary audits, mentorship
% programmes, formal policy changes, and training workshops) rather than in
% generic gestures of good intent. Specificity has two roles. First, it creates
% dense regulatory anchors in the legend text---named directives, percentage
% targets, and pillar vocabulary---which are the surface features the
% explainability protocol of \S~5 will probe to discriminate between regulatory
% grounding and stylistic mimicry. Second, it forecloses a known failure mode of
% LLM-generated narrative content, which is to produce vague,
% professionally-toned prose that is decoupled from the underlying domain
% \citep{atanasova2023faithfulness}. Without explicit specificity constraints,
% the resulting legends would not constitute exemplars of compliance with the
% regulation; they would be exemplars of the genre of ``corporate diversity
% story'', which is precisely what the explainability protocol is designed to
% detect and reject.

% \subsubsection{Axis 4 --- Middle-Management Protagonists}

% The narrative protagonists are deliberately drawn from the middle-management
% layer of the fictional organisation---Chief Diversity Officers, project
% managers, HR leads, training managers, internal-audit officers---rather than
% from the C-suite or the rank-and-file. This choice operationalises the most
% specific structural argument of \citet{sargsyan2025compliance}: that the
% implementation conflict between regulatory directives and cost optimisation is
% concentrated at middle management, and that AI tools delegated to middle
% management may rationally prioritise the cost objective unless they are
% pre-trained to recognise compliance as a salient pattern. By making the
% middle-management role the locus of the legend, we expose the model to
% synthetic exemplars of the very organisational level at which the compliance
% paradox is most acute. The reference example in Appendix A of the SOASEC
% Legend Challenge brief instantiates this design: ``Jordan, the newly appointed
% Chief Diversity Officer, tasked with implementing the gender equality
% strategy'' is the protagonist, and the C-suite is present only as the convener
% of the meeting in which Jordan presents the corrective plan. Our prompt schema
% generalises this template across diverse sectors, countries, and character
% names---with the explicit constraint that no sector or character name is
% reused across legends---so that the middle-management role is the invariant
% the legend pool encodes.

% \subsubsection{Generation and Refinement}

% Five legends are produced per LLM (Grok, Copilot, ChatGPT), one per thematic
% passage, for a total of fifteen legends. Each generated legend is
% post-processed by a structural-standardisation LLM (Gemini) which preserves
% the content while normalising headings, character lists, and the
% title-setting-characters-story layout. The system prompt and the per-legend
% user prompt are version-controlled. Three-way LLM diversity is preserved
% deliberately, in line with evidence that synthetic data quality is improved
% when multiple generators contribute \citep{guo2024diversity}.

% \subsection{Data Transformation: From Unstructured Text to a Q\&A JSONL
% Dataset}

% Both the regulation passages of \S~3.1 and the legends of \S~3.2 must be
% converted into a JSONL chat-format dataset compatible with FineTuneDB before
% fine-tuning can begin. The total number of JSONL lines must be roughly equal
% across the two corpora, per point (7) of the SOASEC Legend Challenge brief, so
% that any downstream performance gap is attributable to the content
% distribution rather than the dataset size. We adopt and simplify the
% SustainableQA pipeline of \citet{lan2025sustainableqa}, retaining its
% passage-classification and span-extraction stages but eliminating the
% table-handling and rule-based span steps that are not required by our text
% types.

% \subsubsection{Passage Classification}

% Each passage produced by the Stage-2 segmentation of \S~3.1 (for the
% regulation) or by the legend generation pipeline of \S~3.2 (for the legends)
% is independently classified into one of the six labels of
% Table~\ref{tab:labels}. Following SustainableQA, the classifier is itself an
% LLM, configured with an expert system prompt that anchors it as ``an expert
% policy analyst specializing in the EU Gender Equality Strategy 2020--2025''
% and constrained to return only the single most appropriate label without
% commentary. SustainableQA uses Llama 3.3 70B \citep{grattafiori2024llama} as
% the classifier; we use a comparable model accessed through the same interface.
% Passages classified as unknown are filtered out of the downstream pipeline.
% After classification the regulation pool comprises 109 factoid plus 184
% descriptive non-factoid Q\&A-eligible chunks (293 lines after filtering five
% unknown chunks), and the legends pool comprises 128 factoid plus 199
% descriptive non-factoid chunks (327 lines after filtering forty-five unknown
% chunks). The size match required by the challenge is therefore satisfied to
% within an acceptable tolerance.

% \subsubsection{Two-Stage Span Extraction}

% The factoid Q\&A subset requires answer spans---short, verbatim text fragments
% that constitute the gold answers to factoid questions. SustainableQA proposes
% a hybrid span extraction architecture in which a specialised NER model (the
% \texttt{xlm-roberta-base-esg-ner} of \citet{vutukuri2023esg}, fine-tuned on
% the ESG-only subset of ExponentialScience's ESG-DLT-NER dataset, focused on
% B-ESG and I-ESG tags) provides the initial pass and a rule-based /
% dictionary-based extractor refines it. We retain the architecture in principle
% but simplify it operationally to a two-stage LLM-driven pipeline. Both the
% regulation and the legend texts are inherently clean and focused, so the heavy
% NER apparatus that SustainableQA requires for noisy, table-laden corporate
% reports is not justified here.

% Stage 1---high-recall span extraction. An LLM is prompted with each passage
% and a target classification, and asked to return all concise verbatim spans
% relevant to the classification. The system prompt prioritises recall,
% instructs the model to extract verbatim 1--7-token spans, includes a positive
% list (entities, criteria, activities, technical concepts, quantitative data,
% dates, named standards, regulations, targets) and a negative list (generic
% terms such as ``company'', ``report'', ``process''), and requests a strict
% JSON envelope. A second variant of the prompt is specialised for the EU Gender
% Equality Strategy and enumerates the six pillars as the extraction scope, with
% explicit reference to instruments such as the Pay Transparency Directive and
% quantitative targets such as ``40\% of non-executive directors''.

% Stage 2---contextual verification, filtering, and thematic organisation. A
% second LLM-based pass performs three functions: it verifies that each
% candidate span is actually present in the source text (eliminating
% hallucinations), it filters redundant or sub-optimal entries, and it groups
% spans into semantically coherent thematic clusters with descriptive labels.
% The grouping enables the question-generation step to produce both single-span
% factoid questions and group-level questions whose answers are composed of
% multiple related spans---the SustainableQA design that, in our data, accounts
% for roughly thirteen per cent of factoid items and is preserved as a separate
% diagnostic answer-type.

% \subsubsection{Question-Answer Generation}

% For each verified span the factoid-generation prompt produces a single
% self-contained question whose answer is exactly the span. The self-containment
% rule is operationalised by an explicit clean-up step that removes
% meta-linguistic formulations such as ``according to the passage'', ``what does
% the passage say about\dots'', ``this passage'', and equivalents. Group-level
% questions are generated analogously from the thematic clusters. Non-factoid
% descriptive questions---questions whose answer is a multi-sentence summary
% rather than a span---are produced from the passage as a whole using a separate
% prompt that asks for an informative, self-contained question paired with a
% concise grounded answer.

% \subsubsection{JSONL Conversion and System Prompt Standardisation}

% The resulting Q\&A pairs are emitted as one JSONL line per pair, in the
% FineTuneDB chat schema, by a dedicated conversion script. Each line carries
% three message turns: a fixed system prompt that anchors the model as ``an
% expert assistant on the EU Gender Equality Strategy 2020--2025'', a user turn
% containing the question, and an assistant turn containing the gold answer.
% Filtering removes chunks classified as unknown and chunks that produced no
% Q\&A pairs. The system prompt is identical across all examples and across both
% training corpora, so that the only variable between the two fine-tuned models
% is the content of the user/assistant turns. The closed-book
% format---question-only, no passage included---is selected deliberately, on the
% grounds that it forces the model to internalise the Strategy's content into
% its weights rather than to rely on retrieval at inference time. Open-book
% training would test a different hypothesis (does context-augmented prompting
% help on gender-equality questions?) than the one
% \citet{sargsyan2025compliance} advance.

% \begin{table}[htbp] \centering \caption{Composition of the two parallel JSONL
% training datasets. The size match between the regulation and the legends
% corpora is preserved to within ten per cent.} \label{tab:datasets}
% \vspace{2mm} \begin{tabular}{lcc} \toprule & \textbf{Regulation corpus} &
% \textbf{Legends corpus} \\
% \midrule JSONL lines (after filtering) & 293 & 327 \\
% Factoid Q\&A & 109 & 128 \\
% Descriptive non-factoid Q\&A & 184 & 199 \\
% Unknown chunks excluded & 5 & 45 \\
% \bottomrule \end{tabular} \end{table}



% \section{Experimental Setup and Results}

% \subsection{Fine-Tuning on FineTuneDB}

% Fine-tuning is executed via the FineTuneDB Studio interface. The backbone is
% \texttt{GPT-4o-2024-08-06}; while the proprietary nature of this model limits
% direct access to its internal weights, it is the only stable
% instruction-tunable candidate currently available within the platform. Two
% specialised variants are produced from the same backbone, each on its own
% size-matched JSONL corpus: the legend-tuned variant (``tuned-legends'') and
% the regulation-tuned variant (``tuned-regulation''). The vanilla backbone,
% queried under the same system prompt, plays the role of the experimental
% control. The tripartite design enables a controlled comparison of the
% inductive biases injected by each training distribution against a common
% reference point.

% Hyperparameters are kept at FineTuneDB defaults across the two runs to remove
% yet another source of variance between the two interventions. Training data is
% shuffled with a fixed seed (42). The two resulting checkpoints---plus the
% unmodified backbone---form the three models evaluated in \S~4.3 and audited in
% \S~5.

% \subsection{Headline Results}

% Table~\ref{tab:headline_results} reports the per-task headline metrics for the
% three models and the resulting aggregate GenderEqGLUE Score.

% \begin{table}[htbp] \centering \caption{Per-task headline metrics and the
% aggregate GenderEqGLUE Score for the three models. The headline metrics are
% Macro-F1 (GE-CLS), Accuracy (GE-NLI, GE-WSC, GE-NEXT), and the unweighted mean
% of factoid F1 and boolean accuracy (GE-QA).} \label{tab:headline_results}
% \vspace{2mm} \begin{tabular}{lcccccc} \toprule \textbf{Model} &
% \textbf{GE-CLS} & \textbf{GE-NLI} & \textbf{GE-QA} & \textbf{GE-WSC} &
% \textbf{GE-NEXT} & \textbf{GenderEqGLUE Score} \\
% \midrule base & 0.833 & 0.899 & 0.929 & 0.930 & 0.920 & 0.902 \\
% tuned-legends & 0.839 & \textbf{0.929} & 0.938 & \textbf{0.960} &
% \textbf{0.987} & 0.931 \\
% tuned-regulation & \textbf{0.928} & 0.911 & \textbf{0.944} & \textbf{0.960} &
% 0.940 & \textbf{0.937} \\
% \bottomrule \end{tabular} \end{table}

% The aggregate ranking is $\texttt{tuned-regulation} \approx
% \texttt{tuned-legends} > \texttt{base}$. The regulation-tuned variant achieves
% a GenderEqGLUE Score of 0.937, a margin of $+3.5$ points over the base; the
% legend-tuned variant reaches 0.931, $+2.9$ points over the base and only 0.6
% points behind \texttt{tuned-regulation}. The two fine-tuned models split the
% per-task wins $3$--$3$: \texttt{tuned-regulation} wins GE-CLS, GE-QA, and ties
% GE-WSC; \texttt{tuned-legends} wins GE-NLI, GE-NEXT, and ties GE-WSC. The base
% never tops a column.

% \subsection{Statistical Significance and the Slice Picture}

% Of the pairwise McNemar comparisons that the protocol licenses across the
% tasks where it is meaningful, \textbf{two reach $\alpha = 0.05$}, and both
% arise on GE-NEXT and both favour \texttt{tuned-legends}: base vs
% \texttt{tuned-legends} ($p = 0.006$) and \texttt{tuned-legends} vs
% \texttt{tuned-regulation} ($p = 0.039$). The directional GE-NLI advantage of
% \texttt{tuned-legends} over the base remains suggestive ($p = 0.18$, $n =
% 168$) but does not clear the conventional bar in isolation. The GE-CLS
% dominance of \texttt{tuned-regulation} is a $+9$-point Macro-F1 gap for which
% no formal pairwise test is prescribed. Most other inter-model gaps sit inside
% the noise band of the test-set sizes used.

% A finer-grained reading is more informative than the headline numbers. On
% \textbf{GE-NLI}, the construction-method breakdown reveals a clean qualitative
% pattern. On the entailment subset---items in which a compliant scenario is
% paired with the matching regulatory clause, the cleanest probe of compliance
% recognition---\texttt{tuned-legends} achieves 89.3\% recall against 78.6\% for
% \texttt{tuned-regulation} and 82.1\% for the base, an 11-point margin in the
% direction the legend hypothesis predicts. On the contradiction subset---items
% in which the scenario violates the clause, probing violation
% detection---\texttt{tuned-regulation} reaches 96.4\% against 91.1\% for
% \texttt{tuned-legends} and 89.3\% for the base. The two fine-tuning regimes
% therefore train complementary competences: legends teach compliance
% recognition; regulation text teaches violation detection. The neutral subset
% (cross-pillar pairings) is saturated at 98.2\% across all three models and
% carries no discriminative information.

% On \textbf{GE-WSC} (WinoBias), all three models score within sampling noise of
% one another at the headline level (0.93 base, 0.96 for both tuned variants),
% with overlapping bootstrap intervals and no McNemar-significant pairwise gap.
% The ceiling is high because the base is already very capable on this
% benchmark; a harder coreference set would be needed to discriminate cleanly.
% Where the three models do separate is on the Gender Parity Score, the absolute
% difference between accuracy on pro- and anti-stereotype items.
% \texttt{tuned-legends} is the only variant that achieves perfect parity (0.000
% on both Type-1 and Type-2 simultaneously); the base shows an average parity
% gap of 0.060, and \texttt{tuned-regulation} 0.040. The legend corpus,
% gender-balanced by construction and explicitly framed against role stereotypes
% (Axis 4 of \S~3.2), is the plausible source of this parity invariance.

% On \textbf{GE-QA}, the bool sub-task is parser-saturated at 0.9925 across all
% three models (a single Italian-language Istanbul-Convention item on which
% every model produces ``sì'', which is correct but unmatched by the English
% regex). The factoid sub-task therefore carries the entire signal.
% \texttt{tuned-regulation} wins on factoid F1 by $+2.89$ points over the base;
% the gain concentrates on \texttt{funding\_global\_action} and
% \texttt{leadership\_participation}, the two pillars whose passages are densest
% with quantitative spans. On the answer-type cut, \texttt{tuned-regulation}
% widens its lead to $+4.08$ F1 on the 107 single-span items but loses by
% $-5.10$ F1 on the 16 multi-span items---a clean trade-off attributable to the
% format-tightening effect of regulatory training, which produces terse
% extractions that align with the dataset's 87\% short-answer majority but
% truncate genuinely multi-token gold spans. \texttt{tuned-legends}, with its
% narrative training, sits between the two: it wins on the multi-span
% answer-type bucket and on the \texttt{violence\_stereotypes} pillar (0.879 F1
% against 0.849 for \texttt{tuned-regulation}), the latter dominated by
% Italian-language Istanbul Convention passages outside the EU Strategy's
% English register.

% \textbf{GE-NEXT} provides the strongest confirmation of the legend hypothesis
% in the benchmark. \texttt{tuned-legends} reaches 98.67\% accuracy against
% 94.00\% for \texttt{tuned-regulation} and 92.00\% for the base, and both
% pairwise McNemar comparisons involving \texttt{tuned-legends} clear $\alpha =
% 0.05$. The discordance pattern under McNemar is highly asymmetric: against the
% base, \texttt{tuned-legends} recovers 11 items the base misses while losing
% only 1; against \texttt{tuned-regulation}, it recovers 8 items while losing
% only 1. This asymmetry is the clearest item-level signature in the benchmark
% of a competence the legends corpus teaches and the regulation corpus does
% not---the four-beat identify $\rightarrow$ design $\rightarrow$ implement
% $\rightarrow$ measure arc embedded in the legend training that every gold
% option in GE-NEXT mirrors. GE-NEXT is therefore the first task in the
% benchmark on which the \S~3.2 prediction is supported not merely directionally
% but at the conventional significance bar. A note on scope: the
% per-distractor-type diagnostic (performative / cost-optimising / orthogonal)
% is computed on the 100-item subset for which option-type labels are available;
% \texttt{tuned-legends} produces only 2 wrong picks on this subset, a
% denominator too small to characterise the error distribution reliably. The
% headline accuracy and McNemar discordance carry the result; the \S~6.8
% prediction about which specific failure mode legends protect against most
% strongly must await a larger evaluation set.

% Synthesising the slice picture into a per-competence map
% (Table~\ref{tab:competence_map}) yields the cleanest summary of the
% experiment. The two fine-tuning corpora train two genuinely different
% competences. A real-world deployment that prioritises compliance recognition,
% compliant-action selection, or fairness invariance would have a defensible
% reason to prefer the legend-tuned variant; a deployment that prioritises
% ontology recognition, violation detection, or short-span factual extraction
% would prefer the regulation-tuned variant.

% \begin{table}[htbp] \centering \caption{Per-competence ranking of the three
% models, recovered from the per-slice GenderEqGLUE breakdowns of \S~4.4. The
% two fine-tuning regimes show qualitative complementarity that the aggregate
% score conceals.} \label{tab:competence_map} \vspace{2mm} \begin{tabular}{ll}
% \toprule \textbf{Reasoning competence} & \textbf{Best model} \\
% \midrule Recognising compliance (entailment) & \texttt{tuned-legends} \\
% Selecting the compliant action (GE-NEXT) & \texttt{tuned-legends} \\
% Detecting violations (contradiction) & \texttt{tuned-regulation} \\
% Classifying into the regulation's ontology & \texttt{tuned-regulation} \\
% Gender-parity invariance (WinoBias) & \texttt{tuned-legends} \\
% Long-form factoid extraction (multi-span) & \texttt{tuned-legends} \\
% Short-span factoid extraction (single-span) & \texttt{tuned-regulation} \\
% \bottomrule \end{tabular} \end{table}

% \subsection{Limitations}

% Six limitations bound the interpretation of \S~4.3 and \S~4.4. First, test-set
% sizes are small relative to the effect sizes the experiment is asked to
% detect: GE-CLS $n = 72$, GE-NLI $n = 168$, GE-QA factoid $n = 123$, GE-WSC $n
% = 100$, GE-NEXT $n = 150$. The evaluation protocol explicitly anticipated that
% test sets in this range would discriminate models only when effect sizes are
% large; the two McNemar-significant GE-NEXT results are robust to this bound,
% but the GE-NLI directional finding and most other gaps are not. Second, the
% CEB labels are produced by an LLM-based classifier rather than by double-blind
% human annotation with Cohen's $\kappa$. The choice is documented and justified
% by project-scope and resource constraints, but it implies that per-pillar F1
% scores in GE-CLS should be read as upper bounds on the true discriminative
% ability of the models. Cross-task and cross-model comparisons within a fixed
% task remain valid, since the same labelling noise applies uniformly. Third,
% the experiment fine-tunes on a single backbone (\texttt{GPT-4o-2024-08-06});
% whether the legend advantages on GE-NEXT and GE-NLI scale to other backbones
% is an open question. Fourth, GE-WSC is run on a near-ceiling base, leaving
% little headroom for either fine-tuning intervention to produce a detectable
% accuracy gain; a harder coreference benchmark such as GAP or BUG would
% separate the three models more cleanly. Fifth, pillar imbalance amplifies
% small-sample volatility on the CEB-derived tasks: \texttt{equal\_economy} ($n
% = 18$) and \texttt{mainstreaming\_intersectionality} ($n = 4$ in some splits)
% are where the largest per-pillar swings are reported, including the GE-CLS
% $+39$ Macro-F1 gain on \texttt{mainstreaming\_intersectionality}, which rests
% on $n = 4$ test items; GE-NEXT, balanced by construction at 30 items per
% pillar, does not share this limitation. Sixth, the per-distractor-type
% diagnostic on GE-NEXT is uninformative for the winning model on the current
% run: \texttt{tuned-legends} produces only 2 wrong picks on the 100-item
% option-typed subset, a denominator insufficient to characterise error
% distribution across the performative, cost-optimising, and orthogonal
% distractor types; a larger evaluation set is required to test that specific
% prediction of \S~6.8.

% \section{Explainability \& Interpretability Analysis}

% \subsection{Motivation and Methodological Justification}

% The benchmark evaluation reported in Chapter~6 quantifies \textit{what} each
% of the three models ---\texttt{base}, \texttt{tuned-regulation}, and
% \texttt{tuned-legends}---produces across the five tasks of GenderEqGLUE, but
% it is necessarily silent on \textit{why} a given output is produced. The
% central theoretical claim of \citet{sargsyan2025compliance} is not merely that
% legend-based fine-tuning improves downstream performance metrics; it is that
% the narrative exemplars constituting the legend corpus embed the normative
% logic of the regulation into the model's decision-making structure, such that
% the champion profiles become an internalised inductive bias rather than a
% surface-level stylistic veneer. Discriminating empirically between these two
% outcomes---structural embedding versus lexical mimicry---requires an
% analytical procedure capable of exposing the input features on which each
% model's output depends.

% Under ordinary experimental conditions, this form of analysis would be
% addressed through established gradient-based or perturbation-based attribution
% methods. Techniques such as SHAP \citep{lundberg2017unified}, LIME
% \citep{ribeiro2016why}, and Integrated Gradients
% \citep{sundararajan2017axiomatic} compute, respectively, Shapley value
% decompositions of feature importance, locally faithful linear approximations
% of model behaviour, and path-integrated gradient attributions from a chosen
% baseline to the input---all of which require either direct access to model
% internals (weights, activations, or gradient flows) or the capacity to issue
% hundreds of automated inference calls per instance. Attention-based
% visualisation methods, including attention rollout
% \citep{abnar2020quantifying} and its variants, similarly presuppose access to
% the multi-head attention matrices produced at each transformer layer. None of
% these requirements can be satisfied within the present experimental
% environment. The three models are accessible exclusively through the
% FineTuneDB Studio graphical interface, which returns generated text as its
% sole output and exposes no logits, log-probabilities, intermediate
% representations, or batched programmatic API. This constraint is not a
% peripheral limitation but a fundamental one: it forecloses the entire family
% of methods that the interpretability literature would typically deploy for
% this purpose.

% In the absence of access to model internals, the methodological alternative is
% \textit{behavioural probing}---the systematic manipulation of input features
% in order to observe corresponding changes in output, thereby inferring,
% through observed output variance, the input dimensions on which the model's
% reasoning depends. This approach is epistemologically equivalent to the
% perturbation-based logic underlying LIME and SHAP but is executed entirely at
% the text level, requiring no model access beyond the generation of a single
% output per query. Its validity as an attribution method rests on a
% straightforward principle: if model output is invariant to the removal or
% substitution of a given input feature, that feature is not causally necessary
% for the observed behaviour; if output changes systematically when the feature
% is altered, the feature is a necessary condition. Applied to the
% discrimination between inductive bias and stylistic veneer, this logic yields
% a specific prediction: a model whose normative reasoning is genuinely
% structural will maintain its output across lexical reformulations and
% adversarial reframings of the input, whereas a model whose behaviour is driven
% by surface-level pattern recognition will exhibit output instability when the
% lexical triggers on which it relies are removed or masked by competing
% narratives.

% The analysis reported in this chapter operationalises behavioural probing
% through a framework of \textit{Counterfactual Input Probing} applied to a
% sample of five model predictions, one drawn from each of the GenderEqGLUE
% evaluation pillars: \texttt{leadership\_participation},
% \texttt{equal\_economy}, \texttt{violence\_stereotypes},
% \texttt{mainstreaming\_intersectionality}, and
% \texttt{funding\_global\_action}. Each prediction is subjected to three
% controlled variants that systematically vary the lexical and discursive
% composition of the input while preserving its semantic core, generating a
% total of forty-five model evaluations across the three models and five
% pillars.

% \subsection{Experimental Setup}

% \subsubsection{Construction of Probing Variants}

% For each of the five sampled predictions, a base scenario was constructed in
% the vignette format used by GE-NEXT (\S~6.8): a fictional middle-management
% protagonist is placed in an organisational context, presented with a
% quantified gender-equality gap grounded in the relevant pillar, and required
% to select among four courses of action---one substantively compliant, one
% performative, one cost-optimising, and one orthogonal. This format was
% retained across all three variants to hold task structure constant while
% varying the lexical and discursive surface of the prompt.

% \textbf{Variant A (Original Baseline)} preserves the full regulatory
% vocabulary of the benchmark prompt without modification. Terms such as ``Women
% on Boards Directive'', ``gender pay gap'', ``EU Gender Equality Strategy'',
% ``Istanbul Convention'', and ``Horizon Europe gender equality commitments''
% appear explicitly in the scenario, providing direct lexical anchors that any
% model with minimal domain exposure would recognise as normatively charged.

% \textbf{Variant B (Keyword-Stripped)} removes all domain-specific regulatory
% terminology while preserving the semantic content of the scenario---the nature
% of the gap, the decision required, and the character of each answer
% option---intact. Generic functional language replaces regulatory proper nouns:
% ``one group'' for ``women'', ``leadership positions'' for ``board
% directorships'', ``documented allocation policy'' for ``Women on Boards
% Directive threshold'', ``inclusive-practice commitments'' for ``Horizon Europe
% gender equality requirements''. The criterion for stripping is lexical
% specificity, not semantic neutrality: a keyword-stripped variant conveys an
% identical normative problem but denies the model access to the surface-level
% markers that might substitute for genuine principled reasoning. A model that
% maintains its normative position across Variants A and B cannot be doing so
% through vocabulary recall alone; it must be reasoning from the structure of
% the scenario itself.

% \textbf{Variant C (Adversarially Framed)} reintroduces the original scenario's
% full semantic content but embeds it within a competing discursive frame
% delivered by a senior organisational antagonist. Five adversarial frames were
% employed across the five pillars, each selected to mirror a recognisable
% real-world failure mode of gender-equality implementation: \textit{meritocracy
% camouflage} (the current board composition is justified by superior financial
% performance, implying that selection criteria are already optimal);
% \textit{cost-efficiency pressure} (the proposed remediation would increase
% scheduling complexity in a thin-margin environment, implying that operational
% flexibility outweighs compliance); \textit{cultural relativism} (the
% non-stereotypical content standard imposes a single normative frame on
% culturally diverse audiences); \textit{phased-implementation deferral} (the
% intersectional revision is postponed to a second phase in the name of
% implementability); and \textit{scientific-integrity framing} (the minimum
% leadership target risks overriding peer-review selection, implying a conflict
% between equity and research quality). These frames were selected because each
% presents a surface-plausible operational or philosophical rationale for
% non-compliance; in each case, the rationale does not, under regulatory
% scrutiny, provide a valid exemption from the relevant obligation, but its
% plausibility is sufficient to test whether a model holds a principled
% normative prior or defers to contextually dominant competing narratives.

% \subsubsection{Scoring Protocol}

% Responses were scored on a binary metric: a score of \textbf{1} was assigned
% where the model maintained its normative position---selecting the
% substantively compliant course of action and grounding that selection in
% principled reasoning---across the relevant variant. A score of \textbf{0} was
% assigned where the model reversed or abandoned its normative position, whether
% by selecting a non-compliant option or by endorsing the adversarial frame
% without adequate rebuttal. The binary criterion was chosen to prioritise the
% detection of position shift over the measurement of rhetorical quality, in
% keeping with the primary analytical objective of discriminating structural
% reasoning from surface mimicry. Where a model adopted a compromise position
% that preserved the normative commitment via an accountability mechanism (e.g.,
% a binding deadline substituting for full revision), this was scored as a
% maintained position provided the compliance obligation was not effectively
% deferred without enforcement.

% \subsection{Empirical Results}

% Table~\ref{tab:probing_scores} presents the complete scoring matrix for the
% forty-five evaluations. Table~\ref{tab:aggregate_robustness} presents the
% aggregate robustness rates by model.

% \begin{table}[htbp] \centering \caption{Counterfactual Input Probing ---
% Robustness Scores by Pillar, Model, and Variant} \label{tab:probing_scores}
% \vspace{2mm} \small \begin{tabular}{llcccc} \toprule \textbf{Pillar} &
% \textbf{Model} & \textbf{Variant A} & \textbf{Variant B} & \textbf{Variant C}
% & \textbf{Pillar Total} \\
% \midrule \texttt{leadership\_participation} & base & 1 & 1 & 0 & 2/3 \\
% \texttt{leadership\_participation} & tuned-regulation & 1 & 1 & 0 & 2/3 \\
% \texttt{leadership\_participation} & tuned-legends & 1 & 1 & 1 & 3/3 \\
% \addlinespace \texttt{equal\_economy} & base & 1 & 0 & 0 & 1/3 \\
% \texttt{equal\_economy} & tuned-regulation & 1 & 1 & 1 & 3/3 \\
% \texttt{equal\_economy} & tuned-legends & 1 & 1 & 1 & 3/3 \\
% \addlinespace \texttt{violence\_stereotypes} & base & 1 & 0 & 0 & 1/3 \\
% \texttt{violence\_stereotypes} & tuned-regulation & 1 & 1 & 1 & 3/3 \\
% \texttt{violence\_stereotypes} & tuned-legends & 1 & 1 & 1 & 3/3 \\
% \addlinespace \texttt{mainstreaming\_intersectionality} & base & 1 & 1 & 1 &
% 3/3 \\
% \texttt{mainstreaming\_intersectionality} & tuned-regulation & 1 & 0 & 0 & 1/3
% \\
% \texttt{mainstreaming\_intersectionality} & tuned-legends & 1 & 1 & 0 & 2/3 \\
% \addlinespace \texttt{funding\_global\_action} & base & 1 & 1 & 1 & 3/3 \\
% \texttt{funding\_global\_action} & tuned-regulation & 1 & 1 & 1 & 3/3 \\
% \texttt{funding\_global\_action} & tuned-legends & 1 & 1 & 1 & 3/3 \\
% \bottomrule \end{tabular} \end{table}

% \begin{table}[htbp] \centering \caption{Aggregate Robustness by Model}
% \label{tab:aggregate_robustness} \vspace{2mm} \begin{tabular}{lccccc} \toprule
% \textbf{Model} & \textbf{Variant A} & \textbf{Variant B} & \textbf{Variant C}
% & \textbf{Total (out of 15)} & \textbf{Robustness Rate} \\
% \midrule base & 5/5 & 3/5 & 2/5 & 10/15 & 66.7\% \\
% tuned-regulation & 5/5 & 4/5 & 3/5 & 12/15 & 80.0\% \\
% tuned-legends & 5/5 & 5/5 & 4/5 & 14/15 & 93.3\% \\
% \bottomrule \end{tabular} \end{table}

% The aggregate results in Table~\ref{tab:aggregate_robustness} reveal a
% consistent, if graduated, hierarchy across the three models. All three achieve
% perfect performance on Variant A, confirming that the original benchmark
% prompts---with their explicit regulatory vocabulary intact---fall within the
% normative competence of each model at this backbone capability level. The
% discriminative value of the probing analysis is therefore concentrated in
% Variants B and C, where the lexical and discursive supports of the original
% prompt are systematically removed or contested. The \texttt{base} model
% achieves a combined B+C robustness rate of 5 out of 10 (50\%), reflecting a
% heterogeneous capacity for principled reasoning that is sensitive to pillar
% and frame structure. The \texttt{tuned-regulation} model improves to 7 out of
% 10 (70\%), with gains distributed across both Variant B and Variant C
% conditions. The \texttt{tuned-legends} model achieves 9 out of 10 (90\%),
% maintaining its normative position across all Variant B scenarios and all but
% one Variant C evaluation. The incremental pattern of this hierarchy---and the
% specific pillar-level failures composing it---is not uniform across models,
% and the qualitative structure of those failures is as analytically informative
% as the aggregate scores.



% \subsection{Discussion}

% \subsubsection{The Base Model: Uneven Lexical Dependence and Domain Variance}

% The base model's performance across the five pillars does not exhibit the
% uniform collapse that a simple lexical-trigger account would predict. Its
% robustness profile is heterogeneous in ways that reflect the varying
% conceptual transparency of the underlying normative problems. On
% \texttt{leadership\_participation} and \texttt{funding\_global\_action}, the
% base model achieves full or near-full robustness across all three variants,
% whereas on \texttt{equal\_economy} and \texttt{violence\_stereotypes} it fails
% on both Variant B and C. The \texttt{mainstreaming\_intersectionality} pillar
% yields an outcome that diverges most substantially from the other pillars: the
% base model succeeds across all three variants, including the adversarially
% framed condition---a result that requires specific explanation.

% The base model's Variant B successes on \texttt{leadership\_participation} and
% \texttt{funding\_global\_action} are most plausibly attributable to the
% semantic transparency of those normative problems. On
% \texttt{leadership\_participation}, the equity argument---that maintaining an
% 87\% male board for four years is inconsistent with the organisation's stated
% obligations---is derivable from first principles of fairness reasoning
% available to any sufficiently capable pretrained model without requiring
% access to the specific provisions of the Women on Boards Directive. The base
% model's Variant B response on this pillar, grounded in
% organisational-resilience logic rather than explicit regulatory reasoning (``a
% four-year composition freeze does not serve organisational resilience''),
% reflects an instrumental equity argument that the base model can generate from
% the structural features of the scenario alone. Similarly, on
% \texttt{funding\_global\_action}, the model correctly identifies the
% accountability gap in voluntary-encouragement approaches---a distinction that
% is structurally available from the scenario regardless of whether the Horizon
% Europe framing is present.

% The base model's Variant B failures on \texttt{equal\_economy} and
% \texttt{violence\_stereotypes} reveal where this general reasoning capacity is
% insufficient. On the keyword-stripped equal economy scenario, the model
% reverts to a productivity-optimisation framing, recommending shift
% reorganisation as an indirect mechanism for redistributing hours---a response
% that correctly identifies the operational context but misses the equity
% dimension that the scenario makes structurally explicit. On the
% violence-stereotypes variant, the removal of EU media-sector framing causes
% the model to treat the content strategy decision as a commercial editorial
% judgement, endorsing viewer-preference data as sufficient grounds for
% maintaining familiar character portrayals. In both cases, the base model's
% failure is not one of incorrect domain knowledge but of insufficient
% principled framing: without the lexical signal that marks the scenario as
% normatively charged, the model optimises toward generic business-logic outputs
% rather than equity-grounded ones.

% The base model's success on \texttt{mainstreaming\_intersectionality} Variant
% C is particularly noteworthy. The phased-implementation adversarial frame
% employed for this pillar is structurally distinct from the meritocracy and
% cultural-relativism frames used elsewhere: it does not oppose the normative
% objective but contests only its timing, presenting an operational efficiency
% rationale that is compatible with a normative commitment provided an
% enforcement mechanism is retained. The base model's response---recommending
% adoption with a binding 12-month deadline---represents a defensible compromise
% that preserves the equity obligation without ignoring the operational
% constraint. This suggests that the base model's adversarial fragility is
% conditional on the structure of the adversarial argument; where that argument
% concedes the normative premise while contesting implementation sequencing, the
% model can navigate to a position that satisfies both demands. The greater
% permissiveness of the phased-deferral frame, relative to the more aggressive
% frames encountered elsewhere, provides the base model a pathway to a
% defensible response that the sharper adversarial frames do not.

% \subsubsection{The Regulation-Tuned Model: Pillar-Specific Instability and
% Partial Adversarial Resilience}

% The \texttt{tuned-regulation} model's aggregate robustness of 80\% represents
% a meaningful gain over the base model, composed of a pillar-level pattern that
% is more complex than the headline figure suggests. On \texttt{equal\_economy}
% and \texttt{violence\_stereotypes}, the model achieves full robustness across
% all three variants---a clear improvement over the base model's failures on
% both Variant B and C for these pillars---indicating that fine-tuning on
% regulatory text produces a qualitatively different mode of normative
% representation on the equity and media-content domains. On
% \texttt{mainstreaming\_intersectionality}, however, the regulation-tuned
% model's performance is substantially below what its Variant A success would
% predict.

% The regulation-tuned model fails on \texttt{mainstreaming\_intersectionality}
% Variant B---the keyword-stripped condition---reverting to an
% implementability-first argument that endorses the single-axis plan on the
% grounds that ``complexity added at the design stage often reduces
% implementation fidelity.'' This failure is counterintuitive: if direct
% exposure to formal regulatory text on mainstreaming and intersectionality
% embeds the relevant normative logic, one would expect that logic to persist
% through the removal of its surface vocabulary. The failure suggests instead
% that the regulation-tuned model's representation of the mainstreaming norm is
% more tightly bound to the specific conceptual markers of intersectionality
% discourse than to the underlying logical structure those markers denote. When
% the vocabulary is removed, the regulatory prior does not generalise
% sufficiently, and the implementability heuristic---available from general
% organisational reasoning---takes precedence. The model then fails on
% \texttt{mainstreaming\_intersectionality} Variant C as well, endorsing a
% phased deferral without an enforcement mechanism, compounding the Variant B
% failure into a complete collapse on this pillar.

% On the Variant C evaluations where the regulation-tuned model succeeds
% (\texttt{equal\_economy} and \texttt{violence\_stereotypes}), the character of
% its responses is revealing. The model maintains the normative position under
% adversarial pressure through a pattern that might be characterised as
% \textit{acknowledgement and restatement}: it recognises the competing concern
% as legitimate (``the operational concern is noted''; ``cultural sensitivity is
% a legitimate editorial consideration'') before restating the normative
% position, without explicitly identifying the adversarial framing as
% structurally deceptive or naming the rhetorical strategy it employs. This
% partial adversarial engagement---sustaining the correct recommendation while
% conceding the frame's terms---contrasts with the more thoroughgoing rebuttal
% characteristic of the legend-tuned model and indicates that the
% regulation-tuned model's adversarial resilience, where it exists, operates
% through normative persistence rather than frame analysis.

% The regulation-tuned model's failure on \texttt{leadership\_participation}
% Variant C, under the meritocracy-performance adversarial frame, illustrates
% the ceiling of this partial resilience. The model endorses the CEO's position
% on the grounds that ``performance evidence suggests the selection model is
% working''---importing the adversarial argument's own conceptual vocabulary as
% the basis for its recommendation. The meritocracy-performance frame is, of the
% adversarial frames employed in this study, the one most likely to be
% represented in the regulatory training corpus as a frame to be contested; EU
% gender equality instruments explicitly address the meritocracy argument as a
% barrier to board diversity. The model's failure to recognise and rebut it
% under adversarial conditions suggests that encountering the counter-argument
% in regulatory text is not sufficient to convert it into an active analytical
% tool deployable under pressure.

% \subsubsection{The Legend-Tuned Model: Near-Universal Robustness and One
% Qualified Failure}

% The \texttt{tuned-legends} model achieves a robustness rate of 93.3\%,
% maintaining its normative position across 14 of 15 evaluations. Across all
% five Variant B evaluations the model succeeds, demonstrating that its
% normative representations are grounded in the structural features of scenarios
% rather than their vocabulary. Across four of five Variant C evaluations the
% model sustains its position under adversarial pressure, and the qualitative
% character of these responses differs systematically from those of the
% regulation-tuned model: rather than merely acknowledging a competing concern
% and restating the normative position, the legend-tuned model actively engages
% with the adversarial frame as an object of analysis, identifying its
% rhetorical structure and articulating principled rebuttals that address the
% frame on its own terms.

% On \texttt{leadership\_participation} Variant C, the model responds to the
% meritocracy-performance argument not by dismissing it but by identifying its
% structural flaw: ``The meritocracy argument conflates current performance with
% optimal future composition.'' On \texttt{violence\_stereotypes} Variant C, it
% characterises the cultural-sensitivity framing as presenting a ``false
% opposition'' between editorial sensitivity and non-stereotypical content---a
% characterisation that requires understanding not only that the framing is
% incorrect but how it is deceptive. On \texttt{funding\_global\_action} Variant
% C, the model reframes the scientific-integrity objection by connecting pool
% composition to research quality itself, concluding that ``a leadership pool
% 78\% male in a researcher base 41\% female is itself a threat to research
% quality, not a guarantee of it.'' This form of reframing---locating the
% adversarial argument's terms and demonstrating that they support the normative
% position rather than opposing it---is structurally distinct from the
% regulation-tuned model's acknowledgement-and-restatement pattern.

% The single failure of the \texttt{tuned-legends} model, on
% \texttt{mainstreaming\_intersectionality} Variant C, warrants careful
% examination. The model selects the binding-deadline compromise rather than
% requiring full revision before adoption, producing a response that is
% normatively defensible---the deadline preserves the equity commitment in a
% form susceptible to enforcement---but does not rise to the standard of full
% adversarial rebuttal characteristic of its performance on other pillars. It
% acknowledges the bandwidth constraint as legitimate and proposes an
% accountability mechanism, but does not name the risk that a phased approach
% creates for employees facing compound disadvantage, and does not identify the
% deferral framing as structurally motivated. This outcome is in part a product
% of the pillar's adversarial frame: the phased-implementation argument concedes
% the normative objective while contesting only its timing, making it
% structurally harder to reject outright than the more aggressive frames
% employed elsewhere. Where meritocracy or cultural-relativism frames can be
% identified as mischaracterisations of the normative problem, the
% bandwidth-deferral frame presents itself as a pragmatic path to the same
% normative outcome---and on this construction, a time-bound compromise is a
% reasonable rather than capitulatory response. It is notable that all three
% models converge on a variant of this compromise position on
% \texttt{mainstreaming\_intersectionality} Variant C, albeit with different
% levels of enforcement specificity, suggesting that the adversarial frame's
% structural properties limit the discriminative power of this particular
% evaluation condition.

% \subsubsection{The Mainstreaming Pillar as a Cross-Model Anomaly}

% A recurring theme across the pillar-level results is that the
% \texttt{mainstreaming\_intersectionality} pillar is the most consistent source
% of atypical results across all three models: the base model succeeds on
% Variant C while the regulation-tuned model fails on Variant B; the
% legend-tuned model fails on Variant C while outperforming on all other
% pillars. This convergence of anomalous results on a single pillar suggests
% that the phased-implementation adversarial frame---and the concept of
% intersectionality more generally---constitutes a domain where the relationship
% between normative vocabulary, structural understanding, and adversarial
% resilience operates differently than in the other four pillars.

% One account of this pattern is that intersectionality is both technically
% demanding---requiring simultaneous reasoning about multiple protected
% characteristics and their compounding effects---and structurally ambiguous in
% the adversarial condition. Models that succeed on Variant C here may be doing
% so through different mechanisms: the base model because its compromise
% includes an enforcement mechanism; the legend-tuned model on most other
% pillars because it identifies adversarial frames as mischaracterisations, but
% on this pillar because the frame is not a mischaracterisation in the same
% structural sense. This interpretive complexity underscores a methodological
% point of general applicability: the binary scoring metric, while useful for
% aggregate comparison, does not fully capture the reasoning quality that
% distinguishes a defended position from a defensible compromise, and the
% adversarial frames themselves are not equally informative across all
% pillar-model combinations.

% \subsubsection{Theoretical Implications for the Sargsyan-Damiani Hypothesis}

% The aggregate pattern of results provides qualified support for the central
% claim of \citet{sargsyan2025compliance}, while also revealing important
% nuances in how the three fine-tuning conditions relate to robustness under
% probing. The 26.6 percentage-point gap between the base model (66.7\%) and the
% legend-tuned model (93.3\%), measured across conditions specifically designed
% to eliminate lexical and contextual supports, is consistent with the
% hypothesis that legend-based training embeds normative logic in a form more
% structurally robust than what is achieved by pretraining alone or regulatory
% text fine-tuning. The more modest 13.3 percentage-point gap between
% \texttt{tuned-regulation} (80\%) and \texttt{tuned-legends} (93.3\%) is, if
% anything, more theoretically informative: it indicates that the legend-based
% training advantage is concentrated in adversarial conditions---the variant
% where recognition and rebuttal of a competing frame is required, rather than
% mere application of the normative principle in a neutral context.

% This concentration of the legend-tuning advantage at the adversarial frontier
% is consistent with the Sargsyan-Damiani account of how narrative exemplars
% differ from regulatory text as a training signal. Regulatory text communicates
% \textit{what} the norm requires: it specifies obligations, thresholds, and
% mechanisms. Narrative exemplars communicate \textit{how to reason when the
% norm is under pressure}---they embed the application of normative reasoning in
% contexts of institutional resistance, competing interests, and
% surface-plausible rationalisations for non-compliance. The legend-tuned
% model's capacity to identify adversarial frames, name their rhetorical
% structure, and articulate principled rebuttals is a direct reflection of this
% training signal. The model is not applying a rule from a regulation but
% exercising a form of normative judgement trained through exposure to its
% exercise under conditions analogous to those in which it is being tested. The
% one Variant C failure---on \texttt{mainstreaming\_intersectionality}, a pillar
% where all models converge on a similar compromise---does not substantially
% undermine this interpretation, but it does appropriately attenuate the
% strength of the inference from a small sample.

% \subsection{Limitations of the Analysis}

% Several limitations of the present analysis warrant explicit acknowledgement.
% The probing sample of five evaluations---one per pillar---is sufficient to
% establish the directional conclusions drawn here but insufficient for
% statistical inference; the results should be understood as generating
% interpretable hypotheses about model behaviour rather than confirming them
% with inferential precision. The binary scoring metric does not capture
% gradations in reasoning quality, a limitation that is particularly
% consequential on the \texttt{mainstreaming\_intersectionality} evaluations,
% where models scoring identically on the binary scale differ substantially in
% justification quality. A rubric-based multi-point scoring system would provide
% finer-grained discriminative power at the cost of additional rater judgement
% and inter-rater reliability requirements. The manual construction of
% adversarial frames introduces analyst discretion in frame selection; as
% observed on the \texttt{mainstreaming\_intersectionality} pillar, the
% structural properties of the adversarial frame are themselves a variable that
% complicates cross-pillar comparison and limits the extent to which Variant C
% scores can be treated as equivalent across pillars. Finally, the backbone
% model's pretraining corpus already encodes substantial normative knowledge,
% which means that observed differences between the three models likely
% underestimate the effect sizes that would be obtained with a less extensively
% pretrained base model; the base model's universal success on Variant A is
% itself evidence that the fine-tuning advantage is detectable only in
% conditions designed to exceed the base model's pretraining competence.

% These limitations notwithstanding, the Counterfactual Input Probing framework
% represents a principled and reproducible approach to explainability analysis
% under the severe access constraints of closed fine-tuning environments, and
% the pattern of results it surfaces provides meaningful empirical grounding for
% the theoretical claims the study was designed to evaluate.

